% Consolidated 23-mechanical-oscillations from 0d2a41587ad18cded191ec1d80e8718df37ec2fe; original TeX fragments preserved.

% === BEGIN 23-mechanical-oscillations-001-020.tex blob b26d3937a8ac14d6118b2e8962da79e4eb1194ec ===
% Irodov Part IV "Oscillations and Waves", §4.1, independent problems 1--20.
% Original uploaded PDF pp.162--164, printed pp.153--155;
% answers PDF pp.322--323, original answers pp.313--314.
\oscproblem{1}{153}
质点沿 $x$ 轴作简谐振动，
$x=a\cos(\omega t-\pi/4)$。
作出(a)位移、速度、加速度随时间变化的示意图；
(b)速度和加速度随坐标 $x$ 变化的图。
\exsolution
求导得
\[
v_x=-a\omega\sin(\omega t-\pi/4),\qquad
a_x=-a\omega^2\cos(\omega t-\pi/4)=-\omega^2x.
\]
三条时间曲线同频率；$v_x$ 比 $x$ 的相位超前 $\pi/2$，$a_x$ 与 $x$ 反相。
消去相位：
\[
\boxed{\left(\frac{x}{a}\right)^2+
\left(\frac{v_x}{a\omega}\right)^2=1},
\qquad
\boxed{a_x=-\omega^2x}.
\]
因此 $(x,v_x)$ 为半轴 $a,a\omega$ 的椭圆，$(x,a_x)$ 为负斜率直线。
\excheck $x=\pm a$ 处 $v_x=0$；$x=0$ 处速度模最大、加速度为零，均与曲线吻合。

\oscproblem{2}{153}
质点沿 $x$ 轴运动，
$x=a\sin^2(\omega t-\pi/4)$。
求(a)振动幅度、周期与 $x(t)$ 图形；
(b)速度 $v_x$ 与坐标 $x$ 的关系及图形。
\exsolution
利用 $\sin^2u=(1-\cos2u)/2$：
\[
x(t)=\frac a2\left[1-\cos(2\omega t-\pi/2)\right],
\]
故它围绕平衡位置 $x_c=a/2$ 振动，
\[
\boxed{A=\frac a2,\qquad T=\frac\pi\omega}.
\]
速度
$v_x=a\omega\sin(2\omega t-\pi/2)$。
因为
$\sin^22u=4\sin^2u(1-\sin^2u)$，
\[
\boxed{v_x^2=4\omega^2x(a-x)},\qquad
\boxed{v_x=\pm2\omega\sqrt{x(a-x)}}.
\]
\excheck 横坐标只能为 $0\le x\le a$；相图闭合在 $x=0,a$，速度模在那里为零。

\oscproblem{3}{154}
质点沿 $x$ 轴作简谐振动，角频率
$\omega=4.00\,\mathrm{s^{-1}}$；
某时刻位置 $x_0=25.0\,\mathrm{cm}$、
速度 $v_0=100\,\mathrm{cm/s}$。
求 $t=2.40\,\mathrm s$ 后的位置与速度。
\exsolution
由运动方程 $\ddot x=-\omega^2x$，
以该时刻为 $t=0$ 的初值解为
\[
x(t)=x_0\cos\omega t+\frac{v_0}{\omega}\sin\omega t,
\quad
v(t)=v_0\cos\omega t-\omega x_0\sin\omega t.
\]
代 $\omega t=9.60$ rad、$x_0=0.25$m、$v_0=1.00$m/s：
\[
\boxed{x(2.40)\simeq-0.29\,\mathrm m},\quad
\boxed{v(2.40)\simeq-0.81\,\mathrm{m/s}}.
\]
\excheck 能量式
$v^2+\omega^2x^2=v_0^2+\omega^2x_0^2=2.0\,\mathrm{m^2/s^2}$
在前后两时刻相等。

\oscproblem{4}{154}
某质点作简谐振动，在距平衡位置分别为
$x_1,x_2$ 时的速度模为 $v_1,v_2$。
求角频率 $\omega$ 和振幅 $a$。
\exsolution
简谐振动能量关系
$v_i^2+\omega^2x_i^2=\omega^2a^2$。
两式相减得到
\[
\boxed{\omega=
\sqrt{\frac{v_1^2-v_2^2}{x_2^2-x_1^2}}}.
\]
代回任一式：
\[
\boxed{a=
\sqrt{\frac{v_1^2x_2^2-v_2^2x_1^2}
{v_1^2-v_2^2}}}.
\]
\excheck 只有上述根号内为正且数据物理相容时才有实数频率和振幅；$a\ge|x_1|,|x_2|$。

\oscproblem{5}{154}
质点简谐振动周期 $T=0.60$s、振幅
$a=10.0$cm。从(a)极端位置；
(b)平衡位置出发，
在走过路程 $a/2$ 的时间内求平均速率。
\exsolution
(a) 从极端位置设 $x=a\cos\omega t$，
走过 $a/2$ 即 $x=a/2$，
于是 $\omega t=\pi/3$，
$\Delta t=T/6$。
平均速率
$\bar v=(a/2)/(T/6)=3a/T$。
(b) 从平衡点设 $x=a\sin\omega t$，
走过 $a/2$ 即 $\sin\omega t=1/2$，
$\Delta t=T/12$，
平均速率 $6a/T$。
故
\[
\boxed{\bar v_a=\frac{3a}{T}=0.50\,\mathrm{m/s}},
\quad
\boxed{\bar v_b=\frac{6a}{T}=1.0\,\mathrm{m/s}}.
\]
\excheck 此区间内运动未反向，因此平均速率也等于位移模除以时间。

\oscproblem{6}{154}
质点从 $t=0$ 开始按
$x=a\sin\omega t$ 运动。
求起始 $3T/8$ 时间内
(a)平均速度的 $x$ 投影；
(b)平均速度矢量的模；
(c)速率的平均值。
\exsolution
$T=2\pi/\omega$，
$\tau=3T/8=3\pi/(4\omega)$，
末坐标
$x(\tau)=a\sin(3\pi/4)=a/\sqrt2$。
故
\[
\boxed{\langle v_x\rangle=
\frac{x(\tau)-x(0)}{\tau}
=\frac{2\sqrt2}{3\pi}a\omega},
\quad
\boxed{|\langle\bm v\rangle|=
\frac{2\sqrt2}{3\pi}a\omega}.
\]
这段时间先从 $0$ 到 $a$，再由 $a$ 折返至 $a/\sqrt2$，
全路程 $s=a+(a-a/\sqrt2)$，
所以
\[
\boxed{\langle|\bm v|\rangle=
\frac{s}{\tau}=
\frac{2(4-\sqrt2)}{3\pi}a\omega}.
\]
\excheck 平均速度与平均速率不同，后者计入折返路程。

\oscproblem{7}{154}
质点按 $x=a\cos\omega t$ 运动，从 $t=0$ 起经过任意时间 $t\ge0$，求累计路程 $s(t)$。
\exsolution
每经历相位 $\pi/2$，质点正好走过距离 $a$。
写 $u=\omega t$，
$n=\lfloor2u/\pi\rfloor$，
$\theta=u-n\pi/2\in[0,\pi/2)$。
完整的 $n$ 个四分之一周期累计路程为 $na$。
剩余相位中，当 $n$ 为偶数，
从端点出发，路程增量 $a(1-\cos\theta)$；
当 $n$ 为奇数，从平衡位置出发，路程增量 $a\sin\theta$。
因此
\[
\boxed{s(t)=
\begin{cases}
a[n+1-\cos\theta],&n\ \text{为偶数},\\
a[n+\sin\theta],&n\ \text{为奇数}.
\end{cases}}
\]
\excheck $t=0$ 时 $n=\theta=0$，得 $s=0$；
每完整周期路程为 $4a$。

\oscproblem{8}{154}
质点自 $t=0$ 起速度投影为
$v_x=35\cos(\pi t)\,\mathrm{cm/s}$，
其中 $t$ 以秒计。
求前 $2.80$s 内走过的路程。
\exsolution
所求路程为绝对值积分
$s=35\int_0^{2.8}|\cos\pi t|\,\mathrm dt$ cm。
速度换向时刻为 $t=0.5,1.5,2.5$s，
逐段积分，得到
\[
s=\frac{35}{\pi}
\bigl[(1)+(2)+(2)+(1-\sin0.8\pi)\bigr]
\,\mathrm{cm}.
\]
即
\[
\boxed{s=\frac{35}{\pi}(6-\sin0.8\pi)
\,\mathrm{cm}\simeq60\,\mathrm{cm}=0.60\,\mathrm m}.
\]
\excheck 不可用最终位移直接代替路程；在 $0.5,1.5,2.5$s 速度依次换向。

\oscproblem{9}{154}
质点按 $x=a\cos\omega t$ 作简谐振动，
假定随机取样时刻在一周期内均匀分布。
求坐标概率密度 $p(x)=\mathrm dP/\mathrm dx$，
其中 $-a<x<a$，并给出图形特征。
\exsolution
对于同一内部位置 $x$，
每周期经过两次，单次停留时间
$\mathrm dt=\mathrm dx/|v_x|$。
因为 $T=2\pi/\omega$，
$v_x^2=\omega^2(a^2-x^2)$，
概率密度
\[
p(x)=\frac{2}{T|v_x|}
=\boxed{\frac1{\pi\sqrt{a^2-x^2}}},\qquad |x|<a.
\]
区间外 $p(x)=0$。
\excheck 归一化
$\int_{-a}^ap(x)\mathrm dx
=\pi^{-1}[\arcsin(x/a)]_{-a}^{a}=1$；
在两端附近密度发散但可积，在 $x=0$ 最小。

\oscproblem{10}{154}
沿同一直线叠加若干同频简谐振动，
用相量法求合振幅：
(a) $x_1=3.0\cos(\omega t+\pi/3)$，
$x_2=8.0\sin(\omega t+\pi/6)$；
(b) $x_1=3.0\cos\omega t$，
$x_2=5.0\cos(\omega t+\pi/4)$，
$x_3=6.0\sin\omega t$。
\exsolution
用复振幅表示 $x=\Re(Ae^{i\omega t})$。
(a) $\sin(\omega t+\pi/6)=\cos(\omega t-\pi/3)$，
所以
$A=3e^{i\pi/3}+8e^{-i\pi/3}$；
相位差 $2\pi/3$，
合振幅
\[
\boxed{|A|=\sqrt{3^2+8^2+2(3)(8)\cos(2\pi/3)}
=\sqrt{49}=7}.
\]
(b) 复振幅
$A=3+5e^{i\pi/4}-6i$。
平方模为
$|A|^2=(3+5/\sqrt2)^2+(5/\sqrt2-6)^2
=70-30/\sqrt2\simeq48.79$，
所以
\[
\boxed{|A|=\sqrt{70-15\sqrt2}\simeq6.99\approx7}.
\]
\excheck (b) 中
$\sin\omega t=\cos(\omega t-\pi/2)$，
因此其相量是 $-6i$ 而不是 $+6i$。

\oscproblem{11}{154}
质点沿同一方向同时参与
$x_1=a\cos\omega t$、
$x_2=a\cos2\omega t$ 两简谐振动。
求合速度的最大模。
\exsolution
总速度
$v=-a\omega[\sin u+2\sin2u]
=-a\omega\sin u(1+4\cos u)$，
$u=\omega t$。
极值条件为
$\mathrm d(\sin u+2\sin2u)/\mathrm du
=\cos u+4\cos2u=0$。
令 $z=\cos u$，
$8z^2+z-4=0$，
有
$z_\pm=(-1\pm\sqrt{129})/16$。
取两候选值及 $u$ 的周期端点，
最大模出现于 $z_+=(-1+\sqrt{129})/16$，
\[
\boxed{v_{\max}=a\omega
\sqrt{1-z_+^2}\,(1+4z_+)
\simeq2.73a\omega}.
\]
\excheck 不能简单将两个分速度振幅
$a\omega$ 与 $2a\omega$ 相加得 $3a\omega$，
因为两个谐波的峰值时刻不同。

\oscproblem{12}{154--155}
质点叠加振动
$x=a\cos(2.1t)\cos(50.0t)$，
$t$ 用秒计。
求两个组成谐波的角频率与拍振周期。
\exsolution
由积化和差：
\[
x=\frac a2[\cos(52.1t)+\cos(47.9t)].
\]
故
\[
\boxed{\omega_1=52.1\,\mathrm{rad/s},
\quad\omega_2=47.9\,\mathrm{rad/s}}.
\]
两频差
$\Delta\omega=4.2\,\mathrm{rad/s}$，
拍振周期
\[
\boxed{T_{\rm beat}=\frac{2\pi}{|\omega_1-\omega_2|}
\simeq1.5\,\mathrm s}.
\]
\excheck 两项等振幅 $a/2$，
包络为 $a|\cos2.1t|$；
相邻包络峰间隔是 $\pi/2.1$ 而非 $2\pi/2.1$。

\oscproblem{13}{155}
质点 A 在参考系 K' 中作某未知频率的简谐振动，
K' 本身沿同一直线相对 K 振动。
当 K' 的频率是 $20$ 或 $24$Hz 时，
在 K 中观察到的拍频均为 $\nu$。
问 K' 振动频率改为多少，拍频变成 $2\nu$？
\exsolution
设 A 相对 K' 的固有振动频率为 $f$。
拍频由两源频率之差给出：
$|f-20|=|f-24|=\nu$，
唯一中点解为
$f=22$Hz、$\nu=2$Hz。
拍频变成 $2\nu=4$Hz 时，
$|f-f'|=4$，故
\[
\boxed{f'=18\,\mathrm{Hz}\ \text{或}\ 26\,\mathrm{Hz}}.
\]
\excheck 两答案都与22Hz差4Hz，
不存在其他非负解。

\oscproblem{14}{155}
平面运动规律
$x=a\sin\omega t$、
$y=b\cos\omega t$，
$a,b,\omega>0$。
求(a)轨迹方程及绕行方向；
(b)加速度矢量关于位置矢量的表达式。
\exsolution
消去时间：
\[
\boxed{\frac{x^2}{a^2}+\frac{y^2}{b^2}=1}.
\]
$t=0$ 时质点在 $(0,b)$，
$v_x=a\omega>0$，
从椭圆顶端向右运动，故沿\emph{顺时针}方向绕行。
分别求二阶导数：
$\ddot x=-\omega^2x$、
$\ddot y=-\omega^2y$，
所以
\[
\boxed{\bm a=-\omega^2\bm r}.
\]
\excheck 即使 $a\ne b$、轨迹为椭圆，
加速度始终指向原点，但速度大小不恒定。

\oscproblem{15}{155}
求平面运动轨迹并作图：
(a) $x=a\sin\omega t$、
$y=a\sin2\omega t$；
(b) $x=a\sin\omega t$、
$y=a\cos2\omega t$。
\exsolution
(a) 令 $u=\omega t$，有
$y=2a\sin u\cos u=2x\cos u$，
而 $\cos^2u=1-x^2/a^2$，
故
\[
\boxed{y^2=4x^2\left(1-\frac{x^2}{a^2}\right)},
\quad |x|\le a.
\]
该轨迹关于两轴对称，是两瓣在原点相交的闭合曲线。
(b) 用
$\cos2u=1-2\sin^2u$：
\[
\boxed{y=a\left(1-\frac{2x^2}{a^2}\right)},
\quad -a\le x\le a.
\]
图形为在有限区间上往复经过的下开口抛物线段。
\excheck (a) 在 $x=0$ 必有 $y=0$，
而 (b) 在 $x=0$ 必有 $y=a$；不可混为同一李萨如图形。

\oscproblem{16}{155}
质量为 $m$ 的质点在一维势场
$U(x)=U_0(1-\cos ax)$ 中，
其中 $U_0,a>0$。
求平衡位置附近的小振动周期。
\exsolution
在稳定平衡点 $x=0$ 附近，
$\cos ax=1-a^2x^2/2+O(x^4)$，
所以
$U(x)\simeq \tfrac12U_0a^2x^2$，
等效弹性系数
$k=U''(0)=U_0a^2$。
角频率
$\omega_0=\sqrt{k/m}=a\sqrt{U_0/m}$，
周期
\[
\boxed{T=\frac{2\pi}{a}\sqrt{\frac m{U_0}}}.
\]
\excheck $a$ 具有长度倒数单位，
所以 $T$ 具有秒的量纲。

\oscproblem{17}{155}
上一题改为
$U(x)=a/x^2-b/x$，其中 $a,b>0$。
求稳定平衡点处小振动周期。
\exsolution
$U'(x)=-2a/x^3+b/x^2$，
稳定点
$x_0=2a/b$。
二阶导数
$U''=6a/x^4-2b/x^3$，
代入
$U''(x_0)=b^4/(8a^3)>0$。
因此
\[
\boxed{T=2\pi\sqrt{\frac{m}{U''(x_0)}}
=\frac{4\pi}{b^2}\sqrt{2ma^3}}.
\]
\excheck 只有正平衡点 $x_0=2a/b$
是本模型内有限稳定点；
二阶导数必须为正才能作小振动。

\oscproblem{18}{155}
一根水平张紧的绳子长
$\ell=1.0\,\mathrm m$，
中点固定一小球，质量 $m=40$g。
绳张力恒定为 $F=10$N。
求小球在竖直平面内的小振动周期。
\exsolution
球竖直位移为 $y$ 时，
每半段绳的长度约为 $\ell/2$，
小偏角
$\sin\theta\simeq y/(\ell/2)=2y/\ell$。
两段绳的回复力
$F_y=-2F\sin\theta\simeq-4Fy/\ell$，
所以
$k_{\rm eff}=4F/\ell$，
\[
\boxed{T=2\pi\sqrt{\frac m{k_{\rm eff}}}
=\pi\sqrt{\frac{m\ell}{F}}
\simeq0.20\,\mathrm s}.
\]
\excheck 原书答案的解析公式与此一致，
但其所印数值“2.0s”与题给40g、10N、1m
直接代入不一致，差十倍；已记录为底本数值疑点。

\oscproblem{19}{155}
绳长 $\ell=20$cm 的单摆，
浸没在密度为摆球密度的 $1/\eta$ 的液体中，
$\eta=3.0$。
忽略流体阻力，求小振动周期。
\exsolution
摆球质量 $m$、体积 $V$，
重力 $mg$，
浮力
$F_b=\rho_\ell Vg=(m/\eta)g$。
有效竖直回复力对应
$g_{\rm eff}=g(1-1/\eta)$，
故
\[
\boxed{T=2\pi\sqrt{
\frac{\ell}{g(1-1/\eta)}}
=2\pi\sqrt{\frac{\eta\ell}{g(\eta-1)}}
\simeq1.1\,\mathrm s}.
\]
\excheck $\eta\to\infty$ 时浮力可忽略，
回到真空中单摆周期；$\eta\downarrow1$ 时周期发散。

\oscproblem{20}{155}
图4.1：长 $\ell$ 的单摆靠近墙壁，静止悬线与墙壁夹小角 $\alpha$；
将摆球再偏离至小角度 $\beta>\alpha$ 后静止释放。
球与墙的碰撞完全弹性。
求其周期。
\exsolution
小角运动满足
$\ddot\theta+\omega_0^2\theta=0$，
$\omega_0=\sqrt{g/\ell}$。
设释放角为 $\theta(0)=\beta$、角速度零，
则
$\theta(t)=\beta\cos\omega_0t$。
按原图的角度约定，碰墙位置是
$\theta=-\alpha$，
到首次碰墙的相位满足
$\cos(\omega_0t_1)=-\alpha/\beta$，
故
$t_1=\omega_0^{-1}[\pi/2+\arcsin(\alpha/\beta)]$。
完全弹性碰撞令角速度反向而不损失机械能，
回到释放端花相同时间，于是
\[
\boxed{T=\frac2{\omega_0}
\left[\frac\pi2+\arcsin\frac\alpha\beta\right]
=2\sqrt{\frac\ell g}
\left[\frac\pi2+\arcsin\frac\alpha\beta\right]}.
\]
\excheck $\alpha\to0$ 时每次仅在摆过竖直位置受墙反弹，周期变为自由单摆周期的一半；
$\alpha\to\beta$ 时碰墙位置靠近另一个转折点，恢复完整简谐周期。
% === END 23-mechanical-oscillations-001-020.tex ===

% === BEGIN 23-mechanical-oscillations-021-033.tex blob 5c2c3293020bb58c83c67c644ba95576547a2351 ===
% Irodov Part IV §4.1 independent original problems 21--33.
% User source PDF pp.165--166; printed pp.156--157.
\oscproblem{21}{156}
单摆钟安置于电梯中。电梯从静止开始以恒加速度
$w$ 向上运动（$w<g$），升高高度 $h$ 时将加速度方向反转而大小仍为 $w$。
问从开始运动起经过多久，钟才重新显示正确时间？
\exsolution
上升加速阶段等效重力
$g_+=g+w$，
摆钟相对于真实时间的走时速率为
$r_+=\sqrt{(g+w)/g}>1$。
从静止上升高度 $h$ 所用真实时间
$t_1=\sqrt{2h/w}$。
反向加速阶段等效重力
$g_-=g-w$，
走时速率
$r_-=\sqrt{(g-w)/g}<1$。
要让累计超前量被追平，
$(r_+-1)t_1+(r_--1)t_2=0$，
所以
$t_2=(r_+-1)t_1/(1-r_-)$。
令 $\eta=w/g$，
\[
\boxed{t=t_1+t_2=
\sqrt{\frac{2h}{w}}\,
\frac{\sqrt{1+\eta}-\sqrt{1-\eta}}
{1-\sqrt{1-\eta}}}.
\]
\excheck 该时刻钟再次准点发生在加速度反向之后；所求为真实经过时间，不是表盘累计时间。

\oscproblem{22}{156}
图4.2密度计（浮计）质量
$m=50$g，细管半径
$r=3.2$mm，浸在液体密度
$\rho=1.00\,\mathrm{g/cm^3}$ 中。
将它沿竖直方向轻微压下后释放，忽略液体阻力，求小振动周期。
\exsolution
下移 $x$ 后额外排开液体体积为
$\Delta V=\pi r^2x$，
增量浮力
$\Delta F_b=\rho g\pi r^2x$，
相对于静平衡的回复力为
$F_x=-\rho g\pi r^2x$。
等效弹性系数
$k=\rho g\pi r^2$，
故
\[
\boxed{T=2\pi\sqrt{\frac{m}{\rho g\pi r^2}}}
=\boxed{2.5\,\mathrm s}\quad\text{（约）}.
\]
\excheck 必须使用浮计\emph{细管}的截面积，
不是浮体整体横截面积。

\oscproblem{23}{156}
一根未变形弹簧两端均固定，
其整体弹性系数为
$k=3.2\,\mathrm{N/m}$。
在距一端为全长
$\eta=1/3$ 的位置附加质量
$m=25$g 的小物体。
忽略重力和弹簧质量，求小物体沿弹簧的微小纵向振动周期。
\exsolution
对同一均匀弹簧材料，
弹簧刚度与其长度成反比。
分段后左段刚度
$k_1=k/\eta$，
右段
$k_2=k/(1-\eta)$。
质量向右移动 $x$ 时，
左段被拉长、右段被压缩，两段均提供回复力，
所以
\[
k_{\rm eff}=k_1+k_2
=\frac{k}{\eta(1-\eta)}.
\]
因此
\[
\boxed{T=2\pi\sqrt{\frac{m\eta(1-\eta)}{k}}
\simeq0.26\,\mathrm s}.
\]
\excheck 本版答案页给出\emph{同一公式}，
但数值标成 $0.13$s；代
$m=0.025$kg、$\eta=1/3$、$k=3.2$N/m
得 $0.26$s，数值相差约2倍，已列勘误。

\oscproblem{24}{156}
图4.3：质量为 $m$ 的小物体在水平光滑轨道上，
左右分别通过刚度为 $k_1,k_2$ 的轻弹簧连接固定支座。
求微小纵向振动周期。
\exsolution
相对平衡位置向右位移 $x$ 时，
左弹簧伸长 $x$、右弹簧缩短 $x$，
两者回复力同向，合力
$F=-(k_1+k_2)x$。
运动方程
$m\ddot x+(k_1+k_2)x=0$，
故
\[
\boxed{T=2\pi\sqrt{\frac{m}{k_1+k_2}}}.
\]
\excheck 两端各连接固定墙的两只弹簧对移动质量而言属于刚度相加，不是串联等效。

\oscproblem{25}{156}
图4.4：一质量为 $m$ 的物体悬挂在\emph{首尾串联}的两根轻弹簧下，
两弹簧刚度分别为 $k_1,k_2$。
求竖直小振动周期。
\exsolution
两弹簧串联时力相同，
总位移 $x=x_1+x_2=F/k_1+F/k_2$，
等效刚度
$k_{\rm eq}=(1/k_1+1/k_2)^{-1}
=k_1k_2/(k_1+k_2)$。
竖直重力只移动静平衡位置，不改变小振动曲率，故
\[
\boxed{T=2\pi\sqrt{\frac{m(k_1+k_2)}{k_1k_2}}}.
\]
\excheck 与上一题不同，图4.4中弹簧确实首尾串联，不能把刚度直接相加。

\oscproblem{26}{156}
质量为 $m$ 的小物体固定在长 $2\ell$ 的水平张紧细绳中点，
绳的静张力为 $T_0$。
忽略重力和绳质量，求横向小振动角频率。
\exsolution
质点离开中点横向位移 $y$，
两侧绳段各长约 $\ell$，
小角度下
$\sin\alpha\simeq y/\ell$，
两侧绳张力的横向合力
$F_y=-2T_0y/\ell$。
所以
$m\ddot y+(2T_0/\ell)y=0$，
\[
\boxed{\omega_0=\sqrt{\frac{2T_0}{m\ell}}}.
\]
\excheck 两段绳各提供一个横向分力，所以刚度带系数2。

\oscproblem{27}{156}
图4.5弯管内装水银总质量
$m=200$g，等截面积
$S=0.50\,\mathrm{cm^2}$，
右臂与竖直夹角
$\theta=30^\circ$。
忽略黏滞，求水银振荡周期。
\exsolution
沿管路令水银柱整体偏移距离 $x$：
左臂液面升高 $x$，
右臂液面沿斜管下降 $x$，
竖直高度改变为 $x\cos\theta$。
两端自由液面竖直高度差为
$(1+\cos\theta)x$，
由液压产生的合回复力
$F=-\rho gS(1+\cos\theta)x$。
整个水银柱质量为 $m$，故
\[
\boxed{T=2\pi
\sqrt{\frac{m}{\rho gS(1+\cos\theta)}}}.
\]
取水银密度 $\rho\simeq13.6\times10^3\,\mathrm{kg/m^3}$，
得
$\boxed{T\simeq0.80\,\mathrm s}$。
\excheck $S$ 的单位换算是
$0.50\,\mathrm{cm^2}=5.0\times10^{-5}\,\mathrm{m^2}$。

\oscproblem{28}{156--157}
图4.6中，一根均匀杆架在两只以相反方向旋转的轮子上。
两轮轴间距 $\ell=20$cm，杆—轮间滑动摩擦系数为
$k=0.18$。
证明杆作简谐振动并求周期。
\exsolution
令杆质心相对两轮中点的位移为 $x$，
左右法向支持力满足
$N_1+N_2=mg$，
力矩平衡
$N_2-N_1=2mgx/\ell$，
所以
$N_1=mg(1/2-x/\ell)$、
$N_2=mg(1/2+x/\ell)$。
图中两轮顶面分别向中心移动，
摩擦力分别为右向 $kN_1$ 与左向 $kN_2$，
合力
$F=k(N_1-N_2)=-2kmgx/\ell$。
因此
\[
\ddot x+\frac{2kg}{\ell}x=0,
\qquad
\boxed{T=2\pi\sqrt{\frac{\ell}{2kg}}
=\pi\sqrt{\frac{2\ell}{kg}}
\simeq1.5\,\mathrm s}.
\]
\excheck 需保持杆始终压在两轮上，
即 $|x|<\ell/2$，此时两支持力均为正。

\oscproblem{29}{157}
设地球为密度均匀的球体，
从北极到南极沿地轴贯穿一条直隧道，
忽略空气阻力和地球自转影响。
求 (a) 质点在隧道中的运动方程；
(b) 从一极到另一极的时间；
(c) 通过地心的速率。
\exsolution
球对称均匀质量分布使距离地心 $x$ 处的引力加速度
$g(x)=-gx/R$，
其中 $R$ 为地球半径、
$g$ 为地表重力加速度。
因此
\[
\boxed{\ddot x+\frac gR x=0}.
\]
频率 $\omega_0=\sqrt{g/R}$，
从一端到另一端是半周期：
\[
\boxed{t_{\rm pole}=\pi\sqrt{\frac Rg}
\simeq42\,\mathrm{min}}.
\]
从地表 $x=R$ 静止释放，地心最大速度
\[
\boxed{v_{\max}=\omega_0R=\sqrt{gR}
\simeq7.9\,\mathrm{km/s}}.
\]
\excheck 原题是沿地球\emph{自转轴}的隧道，
所以忽略自转的离心效应与此几何一致。

\oscproblem{30}{157}
图4.7：长 $\ell$ 的单摆悬点 O 相对地面以恒加速度矢量
$\bm w$ 作任意方向运动，
与重力加速度矢量 $\bm g$ 夹角为 $\beta$。
求小振动周期。计算
$\ell=21$cm、$w=g/2$、$\beta=120^\circ$ 时的周期。
\exsolution
在随悬点加速的非惯性参考系，
有效重力为
$\bm g_{\rm eff}=\bm g-\bm w$，
模
\[
g_{\rm eff}=\sqrt{g^2+w^2-2gw\cos\beta}.
\]
摆球围绕有效重力方向静平衡，小角运动
$\ddot\theta+(g_{\rm eff}/\ell)\theta=0$，
故
\[
\boxed{T=2\pi\sqrt{\frac{\ell}
{\sqrt{g^2+w^2-2gw\cos\beta}}}}.
\]
代 $w=g/2$、$\cos120^\circ=-1/2$，
$g_{\rm eff}=g\sqrt{7/4}$，
得
$\boxed{T\simeq0.80\,\mathrm s}$。
\excheck 不能直接把加速度的模 $w$ 与 $g$ 相减；
需作矢量差。

\oscproblem{31}{157}
图4.8：质量 $m=0.20$kg 的滑套 M 位于两只相同弹簧之间，
两弹簧总有效刚度为 $k=20$N/m，
可沿水平杆 AB 无摩擦滑动。
整个装置绕通过杆中点的竖直轴以恒角速度
$\Omega=4.4\,\mathrm{rad/s}$ 旋转。
求小振动周期，以及使振动不可能出现的 $\Omega$ 范围。
\exsolution
在旋转参考系，
离轴位移 $x$ 的离心力为 $m\Omega^2x$，
两弹簧回复力为 $-kx$，故
\[
m\ddot x=-(k-m\Omega^2)x,
\quad
\omega_{\rm eff}^2=\frac km-\Omega^2.
\]
只在 $k/m>\Omega^2$ 时稳定：
\[
\boxed{T=\frac{2\pi}{\sqrt{k/m-\Omega^2}}
\simeq0.70\,\mathrm s}.
\]
当
$\boxed{|\Omega|\ge\sqrt{k/m}=10\,\mathrm{rad/s}}$
时平衡点不再具有正回复曲率，
不存在该中心平衡附近的小谐振动。
\excheck 阈值等号时频率为零，是临界状态而不是有限周期振动。

\oscproblem{32}{157}
木块置于水平振动木板上，木板作简谐运动，
振幅 $a=10$cm、周期 $T=1.0$s。
当木块刚开始相对木板滑动时，求静摩擦系数的临界值。
\exsolution
木板最大水平加速度为
$a\omega^2$，
$\omega=2\pi/T$。
木块不滑动所需最大水平力为
$ma\omega^2$，
最大静摩擦力为 $\mu mg$。
临界条件
$\mu mg=ma\omega^2$，所以
\[
\boxed{\mu_{\rm crit}=\frac{4\pi^2a}{gT^2}
\simeq0.40}.
\]
\excheck 当 $\mu<\mu_{\rm crit}$ 时开始滑动；
原题英文对幅度与周期的语序有些混乱，
本稿以给出的 $a,T$ 与书末公式为准。

\oscproblem{33}{157}
长 $\ell=80$cm 的单摆在小角条件下，
求三种初始条件对应角位移：
(a) 初偏角 $3.0^\circ$，无初速；
(b) 初在平衡处，摆球获得水平初速 $0.22$m/s；
(c) 原版题面印作初偏角 $30^\circ$，
并获得大小 $0.22$m/s、指向平衡位置的速度。
\exsolution
小角单摆角频率
$\omega_0=\sqrt{g/\ell}=3.5$rad/s，
一般解为
$\theta(t)=\theta_0\cos\omega_0t+
[\dot\theta_0/\omega_0]\sin\omega_0t$。
(a)
\[
\boxed{\theta(t)=3.0^\circ\cos(3.5t)}.
\]
(b) $\dot\theta_0=v_0/\ell=0.275$rad/s，
$|\dot\theta_0|/\omega_0=0.0786$rad$=4.5^\circ$，
所以
\[
\boxed{\theta(t)=4.5^\circ\sin(3.5t)}.
\]
(c) 以印刷题面\emph{字面}初偏角
$30^\circ$ 代入线性化公式得到
\[
\boxed{\theta(t)\approx30^\circ\cos(3.5t)
-4.5^\circ\sin(3.5t)}.
\]
但 $30^\circ$ 已非严格小角，若求精确摆动应解
$\ddot\theta+(g/\ell)\sin\theta=0$。
书末答案却给
$5.4^\circ\cos(3.5t+1.0)$，
它对应初偏角约 $3.0^\circ$
（而非 $30^\circ$），属于明显的原版题面/答案歧义。
\excheck 不把 $30^\circ$ 擅改为 $3.0^\circ$；
书末数据相容性已单独保留，且说明小角近似限制。
% === END 23-mechanical-oscillations-021-033.tex ===

% === BEGIN 23-mechanical-oscillations-034-050.tex blob 27db2a0d9b0706b489782913ac51fc37c91a274c ===
% Irodov Part IV §4.1 original 34--50.
% User uploaded source PDF pp.167--169, printed 158--160; answers pp.324--325.
\oscproblem{34}{158}
图4.9：质量为 $m_1=1.00$kg 的 A 物体与
$m_2=4.10$kg 的 B 物体由一轻弹簧竖直相连，
B 支承在地面上。
A 的自由竖直振动振幅为
$a=1.6$cm、角频率
$\omega=25\,\mathrm{s^{-1}}$。
求装置对地面压力的最大和最小值。
\exsolution
将 A、B 及弹簧整体视为系统，
地面支持力减总重力等于 A 物体的加速度引起的系统竖直合力：
\[
N-(m_1+m_2)g=m_1\ddot y_A.
\]
A 的加速度振幅为 $a\omega^2$，
所以
\[
\boxed{N_{\max,\min}=
(m_1+m_2)g\pm m_1a\omega^2}.
\]
代入
$(m_1+m_2)g\approx50$N、
$m_1a\omega^2=1(0.016)(625)=10$N，
\[
\boxed{N_{\max}\simeq60\,\mathrm N,\quad
N_{\min}\simeq40\,\mathrm N}.
\]
\excheck 两值均为正，说明 B 在本题所用数据下保持与地面接触。

\oscproblem{35}{158}
质量为 $m$ 的小物体放在竖直振动板上，
板自 $t=0$ 起按
$y=a(1-\cos\omega t)$ 上下运动，
$\omega=11\,\mathrm{s^{-1}}$。
求(a) 当 $a=4.0$cm 时物体压板力随时间变化；
(b) 开始脱离板的临界振幅；
(c) 物体在脱板后能上升到相对初始位置
$h=50$cm 所需的振幅。
\exsolution
在未脱离期间，板加速度
$\ddot y=a\omega^2\cos\omega t$，
竖直受力
$N-mg=m\ddot y$。
因此
\[
\boxed{N(t)=mg[1+(a\omega^2/g)\cos\omega t]}.
\]
(a) $a=0.040$m 时
$N/(mg)=1+0.494\cos(11t)$，
全程正值。
(b) 最小压力为 $m(g-a\omega^2)$，
首次可脱离条件 $a\omega^2=g$：
\[
\boxed{a_{\min}=g/\omega^2
\simeq8.1\,\mathrm{cm}}.
\]
(c) 脱板时刻满足
$\cos\omega t_*=-g/(a\omega^2)$。
那时高度
$y_*=a+g/\omega^2$，
向上速度平方
$v_*^2=a^2\omega^2-g^2/\omega^2$。
脱板后再升高度
$v_*^2/(2g)$，
总最高点
$h=a+g/(2\omega^2)+a^2\omega^2/(2g)$，
从而
\[
\boxed{a=\frac{\sqrt{2gh}-g/\omega}{\omega}
\simeq0.20\,\mathrm m}.
\]
\excheck (c) 的振幅大于(b)阈值，确保确实会离板；使用二次抛体运动处理脱离后的上升。

\oscproblem{36}{158}
将质量为 $m$ 的物体悬挂于原先未伸长的竖直弹簧上，
然后无初速度释放；弹簧劲度系数为 $k$，弹簧质量忽略。
求(a)下降位置随时间的关系；
(b)弹簧张力最大与最小值。
\exsolution
以弹簧自然长度处的释放位置为 $y=0$、
向下为正，
运动方程
$m\ddot y=mg-ky$，
令 $\omega=\sqrt{k/m}$，
初值 $y(0)=\dot y(0)=0$，
解为
\[
\boxed{y_{\rm release}(t)=\frac{mg}{k}
[1-\cos(\omega t)]}.
\]
若严格以新静平衡位置 $y_{\rm eq}=mg/k$ 为零，
则
$\boxed{y_{\rm rel.eq}(t)=-(mg/k)\cos\omega t}$。
弹簧张力 $F_s=ky_{\rm release}(t)$，
故
\[
\boxed{F_{s,\max}=2mg,\qquad
F_{s,\min}=0}.
\]
\excheck 原书文字说“相对于平衡位置的位移”，
但书末 $[1-\cos]mg/k$ 明显是相对\emph{释放位置}的下降量；
译稿同时给出两个坐标约定，避免混淆。

\oscproblem{37}{158}
质量为 $m$ 的粒子受中心回复力
$\bm F=-\alpha m\bm r$，
其中 $\alpha>0$。
初始位置 $\bm r(0)=r_0\bm i$，
速度 $\bm v(0)=v_0\bm j$。
求平面运动轨迹。
\exsolution
沿两个互相垂直方向分别满足
$\ddot x+\alpha x=0$、
$\ddot y+\alpha y=0$。
令 $\omega=\sqrt\alpha$，
使用初值：
\[
x(t)=r_0\cos\omega t,\qquad
y(t)=\frac{v_0}{\omega}\sin\omega t.
\]
消去时间：
\[
\boxed{\frac{x^2}{r_0^2}+
\frac{\alpha y^2}{v_0^2}=1}.
\]
\excheck 轨迹是以原点为中心的椭圆；
若 $v_0=\sqrt\alpha\,r_0$，特例为圆。

\oscproblem{38}{158}
质量为 $m$ 的物体悬挂在电梯顶端的轻弹簧下，
劲度系数 $k$。
从 $t=0$ 起电梯向上加速，
物体相对于电梯的初始位移及速度均为零。
求相对于原静平衡位置的运动规律：
(a) 电梯加速度 $w$ 恒定；
(b) $w(t)=\gamma t$。
\exsolution
设 $y$ 是相对电梯原静平衡位置的向下位移。
初始时重力与弹簧力平衡；
在电梯参考系，新增加的惯性力为 $mw(t)$，
所以
\[
\ddot y+\omega_0^2y=w(t),
\qquad \omega_0=\sqrt{k/m},
\quad y(0)=\dot y(0)=0.
\]
(a) 对常数 $w$：
\[
\boxed{y(t)=\frac w{\omega_0^2}
(1-\cos\omega_0t)}.
\]
(b) 用受迫系统格林函数
$y(t)=\omega_0^{-1}\int_0^t
w(s)\sin[\omega_0(t-s)]\,\mathrm ds$，
当 $w=\gamma t$ 时
\[
\boxed{y(t)=\frac{\gamma}{\omega_0^3}
(\omega_0t-\sin\omega_0t)}.
\]
\excheck 两式初始位置、速度为零，
代回微分方程均得到相应驱动加速度。

\oscproblem{39}{158}
质量 $m=0.50$kg 的物体悬挂于弹性系数
$k=50\,\mathrm{N/m}$ 的橡皮绳下。
向下最多拉离平衡位置多少，释放后振动才仍是谐振动？
此临界振动能量是多少？
\exsolution
橡皮绳只在伸长时产生张力，
新静平衡伸长量
$x_0=mg/k$。
若围绕此点振幅 $A>x_0$，
上转折点绳将松弛，回复力不再线性。
临界振幅
\[
\boxed{A_{\max}=\frac{mg}{k}
=\frac{0.50(9.8)}{50}
=9.8\,\mathrm{cm}\approx10\,\mathrm{cm}}.
\]
此时振动能量
\[
\boxed{E_{\max}=\frac12kA_{\max}^2
=\frac{m^2g^2}{2k}\simeq0.24\,\mathrm J}.
\]
\excheck 本版书后解析式与上式一致，
但末尾数值写成约 $4.8\,\mathrm{mJ}$，
与给定 $m,k$ 不符，差约50倍；记录待核验。

\oscproblem{40}{158--159}
图4.10：质量为 $m$ 的小球从静止自由落下高度 $h$，
撞在弹簧秤的无质量托盘上并粘住，
弹簧刚度为 $k$，初始未压缩。
求其后的竖直振动振幅和机械振动能量。
\exsolution
接触前速度
$v_0=\sqrt{2gh}$。
接触后新静平衡弹簧压缩量
$x_0=mg/k$，
相对于新平衡位置初始坐标
$y(0)=-x_0$，初速 $v_0$ 向下。
角频率 $\omega_0=\sqrt{k/m}$，
振幅由初始相空间半径给
\[
A^2=y(0)^2+\frac{v_0^2}{\omega_0^2}
=\left(\frac{mg}{k}\right)^2+\frac{2mgh}{k}.
\]
所以
\[
\boxed{A=\frac{mg}{k}
\sqrt{1+\frac{2hk}{mg}}},
\quad
\boxed{E=\frac12kA^2
=mgh+\frac{m^2g^2}{2k}}.
\]
\excheck 题设托盘无质量，粘住碰撞后球速度仍为
$\sqrt{2gh}$；若托盘有质量则碰撞时损失动能。

\oscproblem{41}{159}
承接第40题，现托盘质量为 $M$，
在物体落下前托盘已停在自身的静平衡位置。
小球质量 $m$ 从高度 $h$ 落下后粘住托盘。
求碰撞后共同振动的振幅。
\exsolution
小球落下速度 $v_0=\sqrt{2gh}$。
完全非弹性粘合碰撞，沿竖直方向短时动量守恒，
新组合体初速
$v_+=m v_0/(m+M)$。
托盘原来压缩 $Mg/k$，
粘合后新平衡压缩
$(M+m)g/k$，
因此相对新平衡初始坐标
$y_0=-mg/k$。
组合质量为 $m+M$，角频率
$\omega_0=\sqrt{k/(m+M)}$，
故
\[
\boxed{A=
\sqrt{\left(\frac{mg}{k}\right)^2+
\frac{m^2(2gh)}{k(m+M)}}
=\frac{mg}{k}
\sqrt{1+\frac{2hk}{(m+M)g}}}.
\]
\excheck $M\to0$ 时精确恢复第40题；碰撞中机械能一般不守恒，必须先守恒动量。

\oscproblem{42}{159}
质量为 $m$ 的粒子在 $xy$ 平面内运动。
原书题面印作速度相关力
$\bm F=\alpha(\dot y\bm i+\dot x\bm j)$，
$\alpha>0$，
初始 $\bm r(0)=0$、$\bm v(0)=v_0\bm j$。
求 $x(t),y(t)$ 及轨迹。
\exsolution
按\emph{原版题面两个正号}，令
$\omega=\alpha/m$，
运动方程为
$\ddot x=\omega\dot y$、
$\ddot y=\omega\dot x$。
初值 $\dot x(0)=0$、$\dot y(0)=v_0$，
解为
$\dot x=v_0\sinh\omega t$、
$\dot y=v_0\cosh\omega t$，
故
\[
\boxed{x(t)=\frac{v_0}{\omega}
(\cosh\omega t-1),\qquad
y(t)=\frac{v_0}{\omega}\sinh\omega t}.
\]
轨迹为
\[
\boxed{\left(1+\frac{\omega x}{v_0}\right)^2-
\left(\frac{\omega y}{v_0}\right)^2=1}.
\]
\emph{书末答案采用另一种力：}
它写
$\ddot y=-\omega\dot x$，
这对应
$\bm F=\alpha(\dot y\bm i-\dot x\bm j)$，
此时才是
$x=(v_0/\omega)(1-\cos\omega t)$、
$y=(v_0/\omega)\sin\omega t$
的圆形轨迹。
\excheck 不能把原题印出的“$+$”悄然改成“$-$”；
本稿保留字面题面的双曲线解及书末另一个符号模型的圆轨迹，待校对印刷源。

\oscproblem{43}{159}
图4.11：轻质薄壁球壳半径 $R$，
内部充满水，球心以轻刚杆悬于点 O，距离为 $\ell$。
水结冰后小振动周期改变多少倍？
忽略水的黏度与结冰体积变化。
\exsolution
结冰前，水在球壳内近似平动而不随壳整体刚性自转，
系统可视为质量为 $m$、有效摆长为 $\ell$ 的单摆，
$J_{\rm liquid}=m\ell^2$。
结冰后水与壳刚性转动，实心冰球绕球心转动惯量
$J_c=(2/5)mR^2$，
绕悬点转动惯量
$J_{\rm ice}=m\ell^2+(2/5)mR^2$。
两态重力回复力矩系数相同为 $mg\ell$，
小摆周期 $T=2\pi\sqrt{J/(mg\ell)}$。
故
\[
\boxed{\frac{T_{\rm ice}}{T_{\rm liquid}}
=\sqrt{1+\frac25\left(\frac R\ell\right)^2}}.
\]
\excheck 比值大于1，结冰后因新增自转惯量而振动变慢。

\oscproblem{44}{159}
图4.12：一根质量 $m$、长 $\ell$ 的均匀细杆，
上端 O 铰接，自然垂直向下；
下端由合刚度为 $k$ 的水平弹簧保持。
求杆的微小平面摆动角频率。
\exsolution
杆相对于 O 转动惯量
$J=m\ell^2/3$。
小转角 $\theta$ 时，重力回复力矩
$\tau_g\simeq-(mg\ell/2)\theta$；
杆底横移 $\ell\theta$，
弹簧回复力 $-k\ell\theta$，
相对 O 的回复力矩
$\tau_s\simeq-k\ell^2\theta$。
所以
\[
J\ddot\theta+
\left(\frac{mg\ell}{2}+k\ell^2\right)\theta=0,
\]
\[
\boxed{\omega_0=
\sqrt{\frac{3g}{2\ell}
\left(1+\frac{2k\ell}{mg}\right)}}.
\]
\excheck $k\to0$ 时还原铰接细杆的物理单摆角频率 $\sqrt{3g/(2\ell)}$。

\oscproblem{45}{159}
图4.13：质量 $m=1.5$kg 的均匀水平杆由两条
长 $\ell=90$cm 的相同细绳悬吊，
绕竖直中心轴转动，
两绳偏离竖直的最大角度
$\alpha=5.0^\circ$，然后释放。
求(a)小转角振动周期；(b)振动能量。
\exsolution
设杆长（两悬点间距）为 $d$，
绕竖直中心轴转动惯量
$J=md^2/12$。
小转角 $\varphi$ 使两绳的横向偏移量约为 $d\varphi/2$，
对应绳倾角 $\alpha\simeq d\varphi/(2\ell)$，
支承高度每条绳缩短
$\Delta h\simeq\ell\alpha^2/2$，
重力势能增加
$U\simeq mg\ell\alpha^2/2
=mgd^2\varphi^2/(8\ell)$。
故扭转回复刚度 $k_\varphi=mgd^2/(4\ell)$，
\[
\boxed{T=2\pi\sqrt{\frac{J}{k_\varphi}}
=2\pi\sqrt{\frac{\ell}{3g}}
\simeq1.1\,\mathrm s}.
\]
最大偏角 $\alpha$ 时全部能量为
\[
\boxed{E=\frac12mg\ell\alpha^2
\simeq0.050\,\mathrm J}.
\]
\excheck $d$ 在周期中消去；$\alpha$ 计算时应换成弧度。

\oscproblem{46}{160}
图4.14：均匀水平圆盘质量为 $m$、半径 $R$，
用扭转系数为 $k$ 的细杆悬挂。
初始扭转角为 $\varphi_0$，
同时给予有向角速度 $\dot\varphi_0$。
求扭转小振动振幅及能量。
\exsolution
圆盘相对中心竖轴的转动惯量
$J=mR^2/2$，
运动方程
$J\ddot\varphi+k\varphi=0$，
频率
$\omega_0=\sqrt{k/J}$。
由初始相空间关系
$\varphi_m^2=\varphi_0^2+
\dot\varphi_0^2/\omega_0^2$，
得到
\[
\boxed{\varphi_m=
\sqrt{\varphi_0^2+\frac{mR^2}{2k}
\dot\varphi_0^{\,2}}},
\quad
\boxed{E=\frac12k\varphi_m^2}.
\]
\excheck 原书对初始角与角速度的文字排印符号接近；
译稿用 $\varphi_0$ 与 $\dot\varphi_0$ 明确区分。

\oscproblem{47}{160}
均匀细杆质量 $m$、长度 $\ell$，
绕上端水平轴作小角振动。
初始偏角为 $\theta_0$、
角速度为 $\dot\theta_0$。
求平均一个周期内的动能。
\exsolution
杆转动惯量
$J=m\ell^2/3$，
重力小角回复系数
$K=mg\ell/2$。
振动总能量
$E=\tfrac12J\dot\theta_0^2+
\tfrac12K\theta_0^2$。
简谐振动中一周期平均动能等于总能量的一半，
所以
\[
\boxed{\langle T_k\rangle=
\frac14J\dot\theta_0^2+
\frac14K\theta_0^2
=\frac{m\ell^2\dot\theta_0^2}{12}
+\frac{mg\ell\theta_0^2}{8}}.
\]
\excheck 小角近似中 $\theta_0$ 须以弧度计；$E/2$ 的等分仅指完整周期平均。

\oscproblem{48}{160}
物理单摆的质心高于悬挂点，
起初处于倒立不稳定平衡附近。
它开始倒下、经过最低稳定平衡位置时的角速度为 $\Omega$。
忽略摩擦，求最低位置附近小角振动周期。
\exsolution
设悬点到质心距离为 $h$，
转动惯量为 $J$。
从上方倒立平衡到下方稳定平衡，重力势能下降
$2mgh$，所以
$\tfrac12J\Omega^2=2mgh$，
即
$mg h/J=\Omega^2/4$。
下方小振动角频率
$\omega_0=\sqrt{mgh/J}=\Omega/2$，
所以
\[
\boxed{T=\frac{2\pi}{\omega_0}
=\frac{4\pi}{\Omega}}.
\]
\excheck 严格在不稳定平衡点零初速会永远不动；
按题意需有任意微小扰动使其开始下落。

\oscproblem{49}{160}
物理单摆原来的小振动角频率
$\omega_1=15.0\,\mathrm{s^{-1}}$；
在距悬点下方 $\ell=20$cm 处固定一小物体，
质量 $m=50$g，
新角频率变成
$\omega_2=10.0\,\mathrm{s^{-1}}$。
求原单摆对悬点的转动惯量。
\exsolution
原摆的重力回复力矩系数设为 $K$，
有
$K=J\omega_1^2$。
增加小质量后
$K'=K+mg\ell$，
$J'=J+m\ell^2$。
新频率满足
\[
\omega_2^2=
\frac{J\omega_1^2+mg\ell}{J+m\ell^2}.
\]
解出
\[
\boxed{J=
\frac{m\ell^2(\omega_2^2-g/\ell)}
{\omega_1^2-\omega_2^2}}
\simeq8.2\times10^{-4}\,\mathrm{kg\,m^2}
=0.82\,\mathrm{g\,m^2}.
\]
\excheck 因 $J>0$，需
$\omega_2^2>g/\ell$；本题数值满足。

\oscproblem{50}{160}
两个物理单摆绕同一水平轴作小振动，
角频率分别为 $\omega_1,\omega_2$，
转动惯量为 $J_1,J_2$。
在稳定平衡位置把它们刚性连在一起，
求复合单摆频率。
\exsolution
单摆小角回复力矩系数
$K_i=J_i\omega_i^2$。
刚性连接后两摆偏角相同，
总转动惯量
$J=J_1+J_2$，
总回复系数
$K=K_1+K_2$。
所以
\[
\boxed{\omega=
\sqrt{\frac{J_1\omega_1^2+J_2\omega_2^2}
{J_1+J_2}}}.
\]
\excheck $\omega^2$ 是两频率平方的惯量加权平均，
因此 $\min(\omega_1,\omega_2)\le\omega\le\max(\omega_1,\omega_2)$。
% === END 23-mechanical-oscillations-034-050.tex ===

% === BEGIN 23-mechanical-oscillations-051-065.tex blob c816769bbffc4dab2f736acf68765bac7edae49e ===
% Irodov Part IV §4.1 original 51--65, source PDF pp.169--172
% printed pp.160--163, figures 4.15--4.23 checked.
\oscproblem{51}{160}
均匀细杆长度 $\ell$，绕穿过杆上某点并垂直于杆的水平轴作小振动。
求悬轴距质心多远时周期最短，并求最小周期。
\exsolution
设轴距质心为 $x>0$，
对轴转动惯量
$J=m(\ell^2/12+x^2)$，
重力回复力矩系数 $K=mgx$。
故
\[
T(x)=2\pi\sqrt{\frac{\ell^2/12+x^2}{gx}}
=2\pi\sqrt{\frac1g\left(x+\frac{\ell^2}{12x}\right)}.
\]
极小点满足
$\mathrm d(x+\ell^2/(12x))/\mathrm dx
=1-\ell^2/(12x^2)=0$，于是
\[
\boxed{x_{\rm opt}=\frac{\ell}{2\sqrt3}},
\qquad
\boxed{T_{\min}=2\pi\sqrt{\frac{\ell}{g\sqrt3}}}.
\]
\excheck 二阶导数
$\ell^2/(6x^3)>0$，
确认是最小而非最大。

\oscproblem{52}{160}
边长相等的均匀薄等边三角形板，
其高度为 $h$。
以其中一条边为水平悬挂轴作小角平面摆动。
求周期和等效摆长。
\exsolution
三角形板的质心距底边
$d=h/3$。
对底边轴转动惯量
$J=m h^2/6$
（可按与底边平行的薄条积分）。
因此
\[
\omega_0^2=\frac{mgd}{J}
=\frac{2g}{h},
\qquad
\boxed{T=2\pi\sqrt{\frac{h}{2g}}
=\pi\sqrt{\frac{2h}{g}}}.
\]
与普通单摆周期
$T=2\pi\sqrt{\ell_{\rm red}/g}$
比较：
\[
\boxed{\ell_{\rm red}=\frac{J}{md}=\frac h2}.
\]
\excheck 对悬轴的转动惯量应为 $mh^2/6$，
不是三角板面内垂直板面的转动惯量。

\oscproblem{53}{161}
图4.15：水平光滑圆盘绕竖直轴 O
以恒角速度 $\Omega$ 转动。
均匀细杆 AB 长 $\ell$，
绕盘上距轴 O 为 $a$ 的竖直轴 A 自由小角摆动。
求相对圆盘径向稳定方向的小振动角频率。
\exsolution
在与圆盘共同转动的参考系，
杆上一质元到 O 的距离平方
$r^2=a^2+2as\cos\theta+s^2$，
其中 $0\le s\le\ell$ 从 A 沿杆测距。
离心势能
$U_{\rm cen}=-\tfrac12\Omega^2\int r^2\,\mathrm dm$，
与 $\theta$ 有关部分
$U_{\rm cen}=-\tfrac12m\Omega^2a\ell\cos\theta$。
小角展开得回复刚度
$K=m\Omega^2a\ell/2$。
杆对 A 的转动惯量
$J=m\ell^2/3$，
因此
\[
\boxed{\omega_0=
\sqrt{\frac KJ}
=\Omega\sqrt{\frac{3a}{2\ell}}}.
\]
\excheck 角频率维度为 $\mathrm{s^{-1}}$，
故不能只写 $\frac{3a\Omega^2}{2\ell}$
（那是角频率的平方）。

\oscproblem{54}{161}
图4.16：半径为 $R$、对轴转动惯量为 $J$ 的滑轮绕轴无摩擦，
一侧悬质量 $m$，另一侧接轻弹簧刚度 $k$。
细线不滑动且无质量。
求平衡位置附近的小振动角频率。
\exsolution
滑轮角位移 $\theta$ 对应悬块及弹簧端点的线位移
$x=R\theta$。
动能
$T=\tfrac12m\dot x^2+
\tfrac12J(\dot x/R)^2
=\tfrac12(m+J/R^2)\dot x^2$。
重力只确定静平衡位置，
微小变形的势能为
$\tfrac12kx^2$。
故
\[
\boxed{\omega_0=
\sqrt{\frac{k}{m+J/R^2}}}.
\]
\excheck 滑轮惯性体现为附加等效质量 $J/R^2$，
不能忽略。

\oscproblem{55}{161}
图4.17：质量为 $M$、半径为 $R$ 的均匀圆柱滑轮，
轴 O 固定且无摩擦；
右侧悬线挂重物 A，
滑轮边缘另一处固定点质量 $m$。
在图示角 $\alpha$ 处静平衡，求小振动角频率。
\exsolution
设悬挂重物 A 的质量为 $m_A$。
平衡时悬重力矩
$m_AgR$ 与边缘点质量的力矩
$mgR\sin\alpha$ 抵消，所以
$m_A=m\sin\alpha$。
小角偏离 $\delta\theta$ 时，
仅边缘点质量的重力矩发生线性回复变化，
刚度为
$K=mgR\cos\alpha$。
总转动惯量
\[
J_{\rm tot}=\frac12MR^2+mR^2+m_AR^2
=R^2\left[\frac M2+m(1+\sin\alpha)\right].
\]
所以
\[
\boxed{\omega_0=
\sqrt{\frac{2mg\cos\alpha}
{R[M+2m(1+\sin\alpha)]}}}.
\]
\excheck 右侧悬重物虽未单独给质量，
但可由平衡条件消去；其平动惯性不可忽略。

\oscproblem{56}{161}
半径为 $r$ 的均匀实心圆柱，
沿半径 $R>r$ 的固定圆柱\emph{内壁}
无滑动滚动，在最低位置附近作小振动。
求周期。
\exsolution
移动圆柱的质心沿半径
$\rho=R-r$ 的圆弧运动。
设质心相对最低点转角为 $\theta$。
无滑动约束使自身角速度模
$\omega_s=\rho\dot\theta/r$。
动能
\[
T_k=\tfrac12m\rho^2\dot\theta^2+
\tfrac12\left(\frac12mr^2\right)
\left(\frac{\rho\dot\theta}{r}\right)^2
=\frac34m\rho^2\dot\theta^2.
\]
势能
$U=mg\rho(1-\cos\theta)
\simeq\tfrac12mg\rho\theta^2$。
故
$\omega_0^2=2g/(3\rho)$，
\[
\boxed{T=2\pi\sqrt{\frac{3(R-r)}{2g}}}.
\]
\excheck 内壁几何决定质心圆轨道半径为 $R-r$，
不能用 $R+r$。

\oscproblem{57}{161}
图4.18：质量为 $m$ 的均匀实心圆柱在水平面上无滑动滚动。
其顶部由左右弹簧牵引，合刚度为 $k$。
求小振动周期。
\exsolution
令圆柱质心水平位移为 $x$，无滑动时转角
$\theta=x/R$。
顶部连接点水平位移
$x+R\theta\simeq2x$，
两弹簧储能
$U=\tfrac12k(2x)^2=2kx^2$，
等效水平刚度 $k_x=4k$。
平动和自转动能
\[
T_k=\frac12m\dot x^2+
\frac12\left(\frac12mR^2\right)
\frac{\dot x^2}{R^2}
=\frac34m\dot x^2,
\]
等效质量 $m_{\rm eff}=3m/2$。
因此
\[
\boxed{T=2\pi\sqrt{\frac{3m/2}{4k}}
=\pi\sqrt{\frac{3m}{2k}}}.
\]
\excheck 图中弹簧连接于圆柱顶端而不是圆心，
因此弹簧端位移约是质心位移的两倍。

\oscproblem{58}{161}
两个质量分别为 $m_1,m_2$ 的立方体，
用刚度 $k$ 的轻弹簧相连，
置于水平光滑面上，
压缩弹簧后同时松手。
求内部相对运动的固有角频率。
\exsolution
以两物体坐标 $x_1,x_2$ 描述，
相对坐标
$\xi=x_2-x_1-\ell_0$，
则
$m_1\ddot x_1=k\xi$、
$m_2\ddot x_2=-k\xi$。
相减：
$\ddot\xi+k(1/m_1+1/m_2)\xi=0$。
定义约化质量
$\mu=m_1m_2/(m_1+m_2)$，
所以
\[
\boxed{\omega_0=\sqrt{k\left(\frac1{m_1}+
\frac1{m_2}\right)}=\sqrt{\frac k\mu}}.
\]
\excheck 质心作匀速运动与弹簧内振动解耦。

\oscproblem{59}{162}
图4.19：两滑球质量
$m_1=1.0$kg、$m_2=2.0$kg，
由刚度 $k=24$N/m 的轻弹簧相连，
在光滑水平杆上滑动。
起初弹簧未伸长，突然给予左球
$v_1=12$cm/s 初速。
求(a)内部振动角频率；
(b)振动能及振幅。
\exsolution
约化质量
$\mu=m_1m_2/(m_1+m_2)=2/3$kg，
相对坐标满足
$\mu\ddot\xi+k\xi=0$，
所以
\[
\boxed{\omega_0=\sqrt{k/\mu}
=\sqrt{24/(2/3)}=6.0\,\mathrm{s^{-1}}}.
\]
起初相对位移 $\xi(0)=0$，
相对速度的模 $|\dot\xi(0)|=v_1$，
振幅
$\boxed{A=v_1/\omega_0=0.020\,\mathrm m=2.0\,\mathrm{cm}}$。
振动能
\[
\boxed{E_{\rm osc}=\frac12\mu v_1^2
\simeq4.8\,\mathrm{mJ}\approx5\,\mathrm{mJ}}.
\]
\excheck 总系统还有质心平动能；
上述 $E_{\rm osc}$ 仅为内部弹簧振动能。

\oscproblem{60}{162}
两个圆盘的轴向转动惯量分别为
$J_1,J_2$，由扭转系数为 $k$ 的细杆连接。
求二者相对小扭转运动周期。
\exsolution
设两圆盘角位移分别
$\theta_1,\theta_2$。
扭转势能
$U=\tfrac12k(\theta_1-\theta_2)^2$。
角动量整体守恒，可用相对角
$\phi=\theta_1-\theta_2$ 描述，
对应约化转动惯量
$J_{\rm red}=J_1J_2/(J_1+J_2)$。
于是
$J_{\rm red}\ddot\phi+k\phi=0$，
\[
\boxed{T=2\pi\sqrt{\frac{J_1J_2}{k(J_1+J_2)}}}.
\]
\excheck 固有整体刚体转动为零频模式；
这里给的是两盘\emph{相对}扭转振动周期。

\oscproblem{61}{162}
图4.20：用两个相同轻弹簧模拟线性
O—C—O 的二氧化碳分子。
三颗小球质量分别为
$m_O,m_C,m_O$，
可按图示两种模态振动。
求两种振动角频率之比。
\exsolution
设每根弹簧刚度 $k$。
对称振型：两端 O 同时向中心或背离中心等幅运动，
中央 C 静止，单端满足
$m_O\ddot x=-kx$，
$\omega_1^2=k/m_O$。
另一振型两 O 同向等幅位移 $x$，
C 反向位移 $y$，总质心守恒
$2m_Ox+m_Cy=0$。
端球的弹力为
$-k(x-y)=-k(1+2m_O/m_C)x$，
故
\[
\omega_2^2=\frac{k}{m_O}
\left(1+\frac{2m_O}{m_C}\right).
\]
因此
\[
\boxed{\frac{\omega_2}{\omega_1}
=\sqrt{1+\frac{2m_O}{m_C}}
=\sqrt{1+\frac{32}{12}}\simeq1.9}.
\]
\excheck 模态顺序依据原版图4.20的箭头；
振动比值与弹簧刚度无关。

\oscproblem{62}{162}
图4.21：充有理想气体的封闭圆筒，
活塞截面积为 $S$，质量为 $m$，
静平衡时将筒内两侧均分成体积 $V_0$，
两侧气压为 $p_0$。
活塞微小位移后释放，
气体变化为绝热过程，忽略摩擦。
求振动角频率。
\exsolution
活塞位移 $x$ 时，
左、右气体体积分别
$V_0+Sx,V_0-Sx$。
由绝热关系 $pV^\gamma=\mathrm{const}$，
一阶展开：
\[
p_L\simeq p_0\left(1-\frac{\gamma Sx}{V_0}\right),
\quad
p_R\simeq p_0\left(1+\frac{\gamma Sx}{V_0}\right).
\]
压力差给回复力
$F=S(p_L-p_R)
=-2\gamma p_0S^2x/V_0$，
所以
\[
\boxed{\omega_0=S
\sqrt{\frac{2\gamma p_0}{mV_0}}}.
\]
\excheck $\gamma=C_p/C_v$ 是绝热指数；
两侧气体都贡献回复力，系数2不可漏。

\oscproblem{63}{162}
图4.22：质量为 $m=21$g 的小球用绝缘细线悬挂在接地无限大水平导体平面上方，
平衡高度 $h=12$cm。
给小球电荷 $q$ 后，振动周期变为原来的 $1/\eta$，
其中 $\eta=2.0$（按书末答案与“changed $\eta$ times”原版语境）。
求 $|q|$。
\exsolution
导体平面镜像等效为距小球 $2h$ 的相反电荷。
球在固定高度附近所受向下静电吸引力
$F_e=q^2/(16\pi\varepsilon_0h^2)$。
在摆动小角度时，球相对平衡点竖直位置变化是二阶小量，
因此可视为有效重力
$g_{\rm eff}=g+F_e/m$，
小振动角频率之比
$\omega_q/\omega_0=\sqrt{g_{\rm eff}/g}=\eta$。
于是
\[
\eta^2=1+\frac{q^2}{16\pi\varepsilon_0h^2mg},
\quad
\boxed{|q|=4h\sqrt{\pi\varepsilon_0mg(\eta^2-1)}
\simeq2.0\,\mathrm{\mu C}}.
\]
\excheck 由于导体吸引加强向下回复力，带电后\emph{周期缩短}；
原版英文仅写“changed $\eta$ times”，
此方向由镜像电力及书后答案共同确认。

\oscproblem{64}{162--163}
小磁针在磁场 $\bm B$ 中绕垂直于磁感应强度的轴作小振动。
改变磁场后周期减小为原来的
$1/\eta$，其中 $\eta=5.0$，
忽略阻尼。
求磁感应强度改变倍数和方向。
\exsolution
磁针磁矩为 $p_m$，
转动惯量 $J$。
小偏角回复力矩
$\tau=-p_mB\theta$，
故
$\omega_0^2=p_mB/J$，
周期
$T=2\pi\sqrt{J/(p_mB)}\propto B^{-1/2}$。
令 $T_f/T_i=1/\eta$，
得到
$\sqrt{B_i/B_f}=1/\eta$，
\[
\boxed{B_f=\eta^2B_i=25B_i}.
\]
\excheck 磁场应\emph{增强}25倍，而不是减小；
摆针本身磁矩与转动惯量保持不变。

\oscproblem{65}{163}
图4.23：两根水平平行导轨间距为 $\ell$，
左端经自感为 $L$ 的理想线圈连接，
右端有质量 $m$ 的活动导杆。
均匀竖直磁场为 $B$，回路总电阻忽略。
$t=0$ 赋予杆向右初速度 $v_0$，
杆初位移为零。
求 $x(t)$。
\exsolution
回路总磁链守恒
$LI+B\ell x=\mathrm{const}$。
以初始电流零、位移零取常数零，
得
$I=-(B\ell/L)x$。
活动杆受安培力
$F_x=IB\ell=-(B^2\ell^2/L)x$。
故
\[
m\ddot x+\frac{B^2\ell^2}{L}x=0,
\qquad
\omega_0=\frac{B\ell}{\sqrt{mL}}.
\]
用初值 $x(0)=0,\ \dot x(0)=v_0$，
\[
\boxed{x(t)=\frac{v_0}{\omega_0}\sin\omega_0t}.
\]
\excheck 无电阻意味着机械能与电感磁场能交换，
没有磁阻尼耗散；轨迹为不衰减正弦振动。
% === END 23-mechanical-oscillations-051-065.tex ===

% === BEGIN 23-mechanical-oscillations-066-081.tex blob ebfd98071e511faa8c86fc3e0241e30a197c32ba ===
% Irodov Part IV §4.1 independent original problems 66--81.
% User PDF pp.172--175, printed pp.163--165, original answers PDF 325--326.
\oscproblem{66}{163}
两根竖直铜轨之间距为 $\ell$，
上端由自感为 $L$ 的线圈短接，
水平活动导杆质量为 $m$，
起初静止，随后沿两导轨自由下落。
整个装置处于垂直导轨平面的匀强磁场 $B$ 中，
各导体电阻忽略。
求导杆下降位移 $x(t)$。
\exsolution
取向下为 $x$ 正向，初始 $x(0)=0$、
$\dot x(0)=0$，初始线圈电流为零。
磁链守恒
$LI+B\ell x=0$，
所以 $I=-(B\ell/L)x$。
活动杆的安培力沿上方、大小
$B^2\ell^2x/L$。
运动方程
\[
m\ddot x=mg-\frac{B^2\ell^2}{L}x.
\]
令
$\omega_0=B\ell/\sqrt{mL}$，
方程化为
$\ddot x+\omega_0^2x=g$。
因此
\[
\boxed{x(t)=\frac{g}{\omega_0^2}
(1-\cos\omega_0t)}.
\]
\excheck 从静止起初加速度为 $g$，
因 $I(0)=0$；回路无电阻时系统能量只在重力势、动能与磁能间交换。

\oscproblem{67}{163}
质点的阻尼振动规律为
$x=a_0e^{-\beta t}\sin\omega t$，
其中 $\beta,\omega>0$。
求 (a) $t=0$ 时位移振幅和速度；
(b) 各个位置转折点所对应的时刻。
\exsolution
幅度包络为
$A(t)=a_0e^{-\beta t}$，
初始包络为
$\boxed{A(0)=a_0}$。
速度
$v=a_0e^{-\beta t}[\omega\cos\omega t-
\beta\sin\omega t]$，
故
$\boxed{v(0)=a_0\omega}$。
转折点满足 $v(t_n)=0$，
即
$\tan\omega t_n=\omega/\beta$，
因此
\[
\boxed{t_n=\frac1\omega
\left[\arctan\frac{\omega}{\beta}+n\pi\right],
\quad n=0,1,2,\ldots}.
\]
\excheck 阻尼会使转折点提前于普通正弦振动的
$\omega t=\pi/2$；$\beta\to0$ 时恢复其极限。

\oscproblem{68}{163}
刚体作阻尼扭转振动，
$\varphi(t)=\varphi_0e^{-\beta t}\cos\omega t$。
求 (a) 初始角速度和角加速度；
(b) 角速度达到局部极值的时刻。
\exsolution
直接求导：
\[
\dot\varphi=-\varphi_0e^{-\beta t}
(\beta\cos\omega t+\omega\sin\omega t),
\]
\[
\ddot\varphi=\varphi_0e^{-\beta t}
[(\beta^2-\omega^2)\cos\omega t+
2\beta\omega\sin\omega t].
\]
所以
\[
\boxed{\dot\varphi(0)=-\beta\varphi_0},
\quad
\boxed{\ddot\varphi(0)=(\beta^2-\omega^2)\varphi_0}.
\]
角速度极值条件 $\ddot\varphi(t)=0$，
\[
\boxed{t_n=\frac1\omega
\left[\arctan\frac{\omega^2-\beta^2}{2\beta\omega}
+n\pi\right]},
\]
其中选取使 $t_n\ge0$ 的整数 $n$。
\excheck 这些是\emph{有符号角速度}的局部极值点；
比较速度模的全部峰值时需区分正峰与负峰。

\oscproblem{69}{163}
阻尼振动写成
$x(t)=a_0e^{-\beta t}\cos(\omega t+\alpha)$。
求初始振幅 $a_0$、初相 $\alpha$，
当初始条件为
(a) $x(0)=0$、$\dot x(0)=v_0$；
(b) $x(0)=x_0$、$\dot x(0)=0$。
\exsolution
一般初值满足
$x_0=a_0\cos\alpha$、
$v_0=-a_0(\beta\cos\alpha+\omega\sin\alpha)$。
(a) $\cos\alpha=0$，
$a_0=|v_0|/\omega$，
并由 $v_0=-a_0\omega\sin\alpha$ 决定相位，
\[
\boxed{a_0=\frac{|v_0|}{\omega},
\quad
\alpha=
\begin{cases}
-\pi/2,&v_0>0,\\
+\pi/2,&v_0<0.
\end{cases}}
\]
(b) $v_0=0$ 给
$\tan\alpha=-\beta/\omega$；
由 $x_0=a_0\cos\alpha$，
\[
\boxed{a_0=|x_0|
\sqrt{1+\left(\frac\beta\omega\right)^2}}.
\]
当 $x_0>0$，
可取
$\alpha=-\arctan(\beta/\omega)$；
当 $x_0<0$，
按相位加 $\pi$ 选择等效分支。
\excheck $\beta=0$ 时退化为标准无阻尼简谐振动初值关系。

\oscproblem{70}{163}
某阻尼振动角频率
$\omega=25\,\mathrm{s^{-1}}$，
在 $t=0$ 时瞬时速度为零，
而实际位移模比该时刻的包络振幅小
$\eta=1.020$ 倍。
求阻尼系数 $\beta$。
\exsolution
写
$x=a_0e^{-\beta t}\cos(\omega t+\alpha)$。
初始速度零给
$\beta\cos\alpha+\omega\sin\alpha=0$，
即
$\tan\alpha=-\beta/\omega$。
位移与包络幅度之比
$|x(0)|/a_0=|\cos\alpha|
=1/\sqrt{1+(\beta/\omega)^2}$。
题设此比为 $1/\eta$，所以
\[
\boxed{\beta=\omega\sqrt{\eta^2-1}
=25\sqrt{1.020^2-1}
\simeq5.0\,\mathrm{s^{-1}}}.
\]
\excheck 对阻尼振动，“速度为零的转折点”
并不恰等于瞬时包络振幅，
其差由 $\beta/\omega$ 决定。

\oscproblem{71}{164}
质点阻尼振动角频率为 $\omega$、
阻尼系数为 $\beta$。
求速度振幅包络随时间的关系，
当 $t=0$：
(a) 位移幅度为 $a_0$；
(b) 实际位移为零，实际速度为 $v_0$。
\exsolution
一般位移
$x=a_0e^{-\beta t}\cos(\omega t+\alpha)$，
速度可写成
\[
v=-a_0e^{-\beta t}\sqrt{\omega^2+\beta^2}
\cos(\omega t+\alpha-\delta)
\]
（适当取相位 $\delta$）。
故
\[
\boxed{V_a(t)=a_0\sqrt{\omega^2+\beta^2}\,
e^{-\beta t}}.
\]
(b) 初始 $x(0)=0$，
按第69题有
$a_0=|v_0|/\omega$，
从而
\[
\boxed{V_b(t)=|v_0|
\sqrt{1+\left(\frac\beta\omega\right)^2}
e^{-\beta t}}.
\]
\excheck 速度包络与位移包络的比例是
$\sqrt{\omega^2+\beta^2}$，不是简单的 $\omega$。

\oscproblem{72}{164}
两种阻尼振动的参数为
$T_1=0.10\,\mathrm{ms}$、
$\beta_1=100\,\mathrm{s^{-1}}$；
$T_2=10\,\mathrm{ms}$、
$\beta_2=10\,\mathrm{s^{-1}}$。
判断哪一种衰减更快。
\exsolution
振幅包络 $A(t)=A_0e^{-\beta t}$。
若比较\emph{相同真实时间}，
$\beta_1>\beta_2$，
所以第一种衰减更快，
其时间常数分别为
$\tau_1=1/\beta_1=0.010$s、
$\tau_2=1/\beta_2=0.10$s。
若比较\emph{每振动一周}衰减程度，
对数减量
$\lambda_i=\beta_iT_i$：
\[
\boxed{\lambda_1=100(10^{-4})=0.010},
\qquad
\boxed{\lambda_2=10(10^{-2})=0.10}.
\]
因此按单周期比例衰减，第二种更快。
\excheck 题目未给“更快”的比较基准，
两个答案分别对应真实时间和自身周期，
与原书书末的提醒一致。

\oscproblem{73}{164}
单摆在某介质中振荡，
对数减量为
$\lambda_0=1.50$。
若介质阻力系数加大
$n=2.00$ 倍，求新对数减量；
阻力加大多少倍后不再出现振荡？
\exsolution
设无阻尼固有角频率为
$\omega_0$，
原阻尼系数 $\beta_0$，
阻尼振动频率
$\omega_{d0}=\sqrt{\omega_0^2-\beta_0^2}$。
对数减量
$\lambda_0=2\pi\beta_0/\omega_{d0}$。
阻力倍数为 $n$ 时
$\beta=n\beta_0$，
$\omega_d=\sqrt{\omega_0^2-n^2\beta_0^2}$，
故
\[
\boxed{\lambda(n)=
\frac{n\lambda_0}
{\sqrt{1+(1-n^2)(\lambda_0/2\pi)^2}}}.
\]
代 $n=2$ 得
$\boxed{\lambda\simeq3.3}$。
临界阻尼为
$n_c\beta_0=\omega_0$，
因此
\[
\boxed{n_c=\sqrt{1+\left(\frac{2\pi}{\lambda_0}\right)^2}
\simeq4.3}.
\]
\excheck $n\ge n_c$ 时系统转为临界或过阻尼，
不再围绕平衡位置反复振荡。

\oscproblem{74}{164}
质量块悬挂在无质量轻弹簧下，
静伸长
$\Delta x=9.8$cm；
轻微竖直推动后阻尼振荡的对数减量为
$\lambda=3.1$。
求实际振荡周期。
\exsolution
静平衡
$k\Delta x=mg$，
无阻尼固有角频率
$\omega_0^2=k/m=g/\Delta x$。
对于欠阻尼振荡，
$\omega_d^2=\omega_0^2-\beta^2$，
周期 $T=2\pi/\omega_d$，
对数减量
$\lambda=\beta T$。
消去 $\beta$：
\[
\frac{4\pi^2}{T^2}
=\frac g{\Delta x}-\frac{\lambda^2}{T^2},
\]
故
\[
\boxed{T=
\sqrt{\frac{(4\pi^2+\lambda^2)\Delta x}{g}}
\simeq0.70\,\mathrm s}.
\]
\excheck 若错误采用无阻尼周期
$2\pi\sqrt{\Delta x/g}$，
将忽略 $\lambda=3.1$ 的显著阻尼效应。

\oscproblem{75}{164}
某振子经过
$n=110$ 次振动后，位移幅度降低
$\eta=2.0$ 倍。
求品质因数 $Q$。
\exsolution
若每周期振幅比
$A_{k+1}/A_k=e^{-\lambda}$，
则经过 $n$ 周期有
$A_n/A_0=e^{-n\lambda}=1/\eta$，
所以
$\lambda=(\ln\eta)/n$。
欠阻尼振子的品质因数按
$Q\simeq\pi/\lambda$，故
\[
\boxed{Q\simeq\frac{\pi n}{\ln\eta}
=\frac{\pi(110)}{\ln2}
\simeq5.0\times10^2}.
\]
\excheck 这里 $\lambda\ll1$，
因此用 $Q\simeq\pi/\lambda$ 的弱阻尼近似非常好。

\oscproblem{76}{164}
粒子初始离开平衡位置 $\ell=1.0$cm，
随后由静止释放，
发生阻尼振动直至停止。
对数减量
$\lambda=0.020$。
求整个过程中累计走过的总路程。
\exsolution
相邻两个转折点间隔半个阻尼周期，
对应位移极值的模按因子
$r=e^{-\lambda/2}$ 递减。
起始极值为 $\ell$，
之后转折点依次为
$-\ell r,+\ell r^2,\ldots$。
总路径为
\[
s=\ell+\ell r+
(\ell r+\ell r^2)+(\ell r^2+\ell r^3)+\cdots
=\ell\frac{1+r}{1-r}.
\]
因此
\[
\boxed{s=\ell
\frac{1+e^{-\lambda/2}}{1-e^{-\lambda/2}}
\simeq2.0\,\mathrm m}.
\]
\excheck 虽然每次反转位移很小，
衰减率仅每半周期约1%，
累计路程可远大于初始1cm振幅。

\oscproblem{77}{164}
数学单摆长度 $\ell=50$cm，
在
$\tau=5.2$min 内总机械能下降
$\eta=4.0\times10^4$ 倍。
求品质因数。
\exsolution
阻尼振动能量包络
$W(t)\propto e^{-2\beta t}$，
故
$\eta=e^{2\beta\tau}$，
$\beta=(\ln\eta)/(2\tau)$。
无阻尼固有角频率
$\omega_0=\sqrt{g/\ell}$，
阻尼角频率
$\omega_d=\sqrt{\omega_0^2-\beta^2}$。
用本书弱阻尼定义
$Q=\omega_d/(2\beta)$，
\[
\boxed{Q=\frac12
\sqrt{\frac{4g\tau^2}{\ell(\ln\eta)^2}-1}
\simeq1.3\times10^2}.
\]
\excheck $5.2$min 需换为312s；
$Q\gg1$ 支持弱阻尼模型。

\oscproblem{78}{164}
均匀薄圆盘半径
$R=13$cm，
可绕垂直盘面、穿过盘缘的水平轴转动，
作小幅阻尼摆动，对数减量
$\lambda=1.00$。
求振荡周期。
\exsolution
圆盘对悬轴的转动惯量
$J=J_{\rm cm}+mR^2=(3/2)mR^2$。
质心距悬轴 $R$，
故无阻尼角频率
$\omega_0^2=mgR/J=2g/(3R)$。
对数减量 $\lambda=\beta T$，
阻尼角频率
$\omega_d=2\pi/T$，
于是
$\omega_0^2=(4\pi^2+\lambda^2)/T^2$，
\[
\boxed{T=
\sqrt{\frac{3R(4\pi^2+\lambda^2)}{2g}}
\simeq0.90\,\mathrm s}.
\]
\excheck 悬轴在\emph{盘缘}且轴垂直盘面，
所以 $J$ 不是圆盘中心惯量 $mR^2/2$。

\oscproblem{79}{164}
水平薄均匀圆盘质量为 $m$、半径为 $R$，
由弹性细线悬挂，扭转回复力矩大小
$\alpha|\varphi|$。
浸于黏性流体中，
每单位盘面面积的阻力大小
$\eta v$（$v$ 为局部速度）。
求小扭转振动角频率。
\exsolution
圆盘绕轴转动惯量
$J=mR^2/2$。
半径 $r$ 处的速度
$v=r|\dot\varphi|$，
每个表面元受到切向阻力密度
$\eta r|\dot\varphi|$，
相对轴的力矩元
$\eta r^2\dot\varphi\,\mathrm dS$。
盘上下两面总阻力矩系数
\[
b=2\eta\int_{\rm disc}r^2\,\mathrm dS
=2\eta\left(\frac{\pi R^4}{2}\right)
=\pi\eta R^4.
\]
方程
$J\ddot\varphi+b\dot\varphi+\alpha\varphi=0$。
设阻尼系数
$\beta=b/(2J)=\pi\eta R^2/m$，
则
\[
\boxed{\omega=
\sqrt{\frac{2\alpha}{mR^2}-
\left(\frac{\pi\eta R^2}{m}\right)^2}}.
\]
\excheck 只有根号内为正时存在欠阻尼振动；
两面流体黏性阻力均需计入。

\oscproblem{80}{165}
图4.24：半径为 $R$ 的薄盘 A
在上下两个固定平面之间扭转振动，
盘对中心轴转动惯量为 $J$，
盘与每个平面的间隙为
$h\ll R$。
振动周期 $T$、对数减量 $\lambda$ 已知。
求周围气体的动力黏度 $\eta$。
\exsolution
半径 $r$ 处盘面切向速度
$v=r\dot\varphi$，
在狭缝内近似线性速度剖面，
切应力为
$\tau_s=\eta v/h=\eta r\dot\varphi/h$。
上下两面总阻力矩
\[
M_{\rm drag}=-2\frac{\eta\dot\varphi}{h}
\int_{\rm disc}r^2\,\mathrm dS
=-\frac{\pi\eta R^4}{h}\dot\varphi.
\]
即线性阻尼系数
$b=\pi\eta R^4/h$。
运动方程
$J\ddot\varphi+b\dot\varphi+k\varphi=0$，
位移包络衰减率
$\beta=b/(2J)$，
对数减量 $\lambda=\beta T$。
所以
\[
\boxed{\eta=\frac{2\lambda hJ}{\pi R^4T}}.
\]
\excheck $\eta$ 维度为
$\mathrm{Pa\,s}$；
本式要求层流薄隙 $h\ll R$ 以及线性黏滞近似。

\oscproblem{81}{165}
图4.25：边长为 $a$ 的正方形导线框
悬在弹性细线下，
处于均匀水平磁场 $B$ 中。
平衡时框面平行磁场，
绕通过框中心的竖轴作微小扭转振动，
框对轴惯量为 $J$、电阻为 $R$，
忽略自感。
求振动角幅度衰减到原来的
$1/e$ 所需时间。
\exsolution
微小转角 $\varphi$ 时，
框面法向磁通
$\Phi=Ba^2\sin\varphi
\simeq Ba^2\varphi$。
感生电动势
$\mathcal E=-Ba^2\dot\varphi$，
感生电流
$I_{\rm ind}=\mathcal E/R$。
磁力矩
$\tau_m=I_{\rm ind}\,\mathrm d\Phi/\mathrm d\varphi
\simeq-(B^2a^4/R)\dot\varphi$，
等效黏性阻尼系数
$b=B^2a^4/R$。
微小扭转方程
$J\ddot\varphi+b\dot\varphi+k\varphi=0$，
在欠阻尼条件下振幅包络为
$A(t)=A_0e^{-\beta t}$，
$\beta=b/(2J)$。
所以
\[
\boxed{\tau_{1/e}=\frac1\beta
=\frac{2JR}{a^4B^2}}.
\]
\excheck 题设未给扭转弹性常数 $k$，
但幅度包络的 e 倍衰减时间只由
惯量和线性电磁阻尼系数决定。
% === END 23-mechanical-oscillations-066-081.tex ===

% === BEGIN 23-mechanical-oscillations-082-093.tex blob 7c03de5b62cfb884700677ecd1e66d0b5df51d96 ===
% Irodov Part IV §4.1 original 82--93.
% Checked user-uploaded Arihant PDF pp.174--176, printed pp.165--167,
% and book answers PDF pp.325--326 (printed pp.316--317).
\oscproblem{82}{165}
质量 $m=0.50\,\mathrm{kg}$ 的木块在水平面上由原长弹簧连接墙面；滑动摩擦系数 $\mu=0.10$，弹簧刚度 $k=2.45\,\mathrm{N/cm}$，弹簧质量忽略。先将弹簧拉长 $x_0=3.0\,\mathrm{cm}$，随后释放。求：(a) 振动周期；(b) 完全停下以前的振动次数。
\exsolution
块在每段单向滑行中受到大小为 $\mu mg$ 的恒定摩擦力，运动方程
$m\ddot x+kx=-\mu mg\,\operatorname{sgn}\dot x$。
分段仍是角频率 $\omega_0=\sqrt{k/m}$ 的简谐运动，因此
\[
\boxed{T=2\pi\sqrt{\frac mk}
=2\pi\sqrt{\frac{0.50}{245}}\simeq0.28\,\mathrm s}.
\]
记 $\Delta=\mu mg/k=0.0020\,\mathrm m=0.20\,\mathrm{cm}$。
从一个转折点运动到下一反向转折点时，由摩擦功或围绕偏移平衡点求解，幅值下降 $2\Delta$；每一整周下降 $4\Delta$。到最后一个转折点，若剩余弹簧力不超过最大静摩擦力（即 $|x|\le\Delta$），物块停住。按原题 $x_0/\Delta=15$，恰经过七个半周后到达 $|x|=\Delta$，所以
\[
\boxed{N=\frac{x_0-\Delta}{4\Delta}
=\frac{3.0-0.20}{0.80}=3.5\ \text{周}}.
\]
\excheck 源题摩擦系数原记为 $k$、弹簧系数原记为 $\varkappa$；本稿用 $\mu,k$ 区别。临界转折点处所需静摩擦不超过 $\mu mg$。

\oscproblem{83}{165}
质量为 $m$ 的无阻尼谐振子固有角频率 $\omega_0$，$t=0$ 时处于平衡位置且静止；从此时起施加外力 $F_x=F_0\cos\omega t$。求随后位移 $x(t)$。
\exsolution
微分方程
$\ddot x+\omega_0^2x=(F_0/m)\cos\omega t$，
初值 $x(0)=\dot x(0)=0$。
当 $\omega\ne\omega_0$，特解与齐次解叠加并由初值得
\[
\boxed{x(t)=\frac{F_0/m}{\omega_0^2-\omega^2}
\left(\cos\omega t-\cos\omega_0t\right)}.
\]
利用 $\cos a-\cos b$ 的极限求共振：
\[
\boxed{x_{\rm res}(t)=\frac{F_0}{2m\omega_0}
\,t\sin\omega_0t}\quad(\omega=\omega_0).
\]
\excheck 非共振时初始位移和速度均为零；共振时振幅包络随 $t$ 线性增长，只有在无阻尼理想模型中成立。

\oscproblem{84}{165}
质量 $m$ 的一维粒子在回复力 $-kx$ 下作自由简谐振动。原先位于平衡点静止，突然施加恒力 $F$ 并持续时间 $\tau$，然后撤去。求撤力后的振动振幅，讨论特殊持续时间并描述位移曲线。
\exsolution
令 $\omega_0=\sqrt{k/m}$。施力期间
$\ddot x+\omega_0^2x=F/m$，
零初值解
\[
x(t)=\frac Fk(1-\cos\omega_0t),\quad
\dot x(t)=\frac{F\omega_0}{k}\sin\omega_0t
\quad(0\le t\le\tau).
\]
撤力后的简谐振幅由撤力瞬间的相空间位置决定：
\[
a=\sqrt{x(\tau)^2+\frac{\dot x(\tau)^2}{\omega_0^2}}
=\frac{|F|}{k}\sqrt{(1-\cos u)^2+\sin^2u},
\quad u=\omega_0\tau.
\]
故
\[
\boxed{a=\frac{2|F|}{k}
\left|\sin\frac{\omega_0\tau}{2}\right|}.
\]
若 $\tau=2\pi n/\omega_0$，撤力时返回原平衡点并静止，$a=0$；若 $\tau=(2n+1)\pi/\omega_0$，$a=2|F|/k$ 达到最大。
\excheck 两段位移及速度在 $t=\tau$ 连续，只有加速度因外力撤去而跳变。

\oscproblem{85}{165--166}
质量为 $m$ 的弹簧振子悬挂后使弹簧伸长 $\Delta\ell$。受振幅为 $F_0$ 的竖直谐波外力，阻尼振动的对数减量为 $\lambda$。弹簧质量忽略，求位移响应幅度最大时的外力角频率和最大位移幅度。
\exsolution
由静平衡 $k\Delta\ell=mg$，固有角频率
$\omega_0=\sqrt{g/\Delta\ell}$。
令阻尼系数 $\beta$，阻尼自由振动频率
$\omega_d=\sqrt{\omega_0^2-\beta^2}$，
对数减量
$\lambda=2\pi\beta/\omega_d$。
记 $\delta=\lambda/(2\pi)$，
则
$\beta^2=\omega_0^2\delta^2/(1+\delta^2)$。
稳态位移幅度
\[
A(\omega)=\frac{F_0/m}
{\sqrt{(\omega_0^2-\omega^2)^2+4\beta^2\omega^2}}.
\]
对分母平方求极小得到
$\omega_{\rm res}^2=\omega_0^2-2\beta^2$，因而
\[
\boxed{\omega_{\rm res}=
\sqrt{\frac{g}{\Delta\ell}}
\sqrt{\frac{1-\delta^2}{1+\delta^2}}},
\qquad \delta=\frac{\lambda}{2\pi}.
\]
在 $\delta<1$ 且频率非零的共振区间，代回分母：
\[
\boxed{A_{\rm res}
=\frac{F_0\Delta\ell}{4\pi mg}
\left(\lambda+\frac{4\pi^2}{\lambda}\right)}.
\]
\excheck 当 $\delta\ge1$ 时非零共振峰消失，幅度最大转为静态极限 $\omega=0$，上式不能无条件继续应用。

\oscproblem{86}{166}
某受迫阻尼振子在外力角频率 $\omega_1=400\,\mathrm{s^{-1}}$ 和
$\omega_2=600\,\mathrm{s^{-1}}$ 时有相同位移振幅。求位移幅度共振角频率。
\exsolution
位移响应幅度分母的平方
$D(u)=(\omega_0^2-u)^2+4\beta^2u$，
$u=\omega^2$。
因 $D(\omega_1^2)=D(\omega_2^2)$ 且两频不同，二次多项式对称轴为
\[
u_{\rm res}=\omega_0^2-2\beta^2
=\frac{\omega_1^2+\omega_2^2}{2}.
\]
于是
\[
\boxed{\omega_{\rm res}=
\sqrt{\frac{\omega_1^2+\omega_2^2}{2}}
=\sqrt{260000}\,\mathrm{s^{-1}}
\simeq5.1\times10^2\,\mathrm{s^{-1}}}.
\]
\excheck 共振频率不是两频率的算术平均 $500\,\mathrm{s^{-1}}$，而是频率平方的平均值开根号。

\oscproblem{87}{166}
谐波外力频率分别为 $\omega_1,\omega_2$ 时，阻尼受迫振子的\emph{速度幅度}都恰为全频段最大速度幅度的一半。求：(a) 速度共振角频率；(b) 阻尼系数 $\beta$ 与自由阻尼角频率 $\omega_d$。
\exsolution
速度响应幅度
$V(\Omega)=(F_0/m)\Omega/
\sqrt{(\omega_0^2-\Omega^2)^2+4\beta^2\Omega^2}$。
其最大值在 $\Omega=\omega_0$，为 $F_0/(2m\beta)$。
半峰条件 $V(\Omega)=V_{\max}/2$ 得
\[
(\omega_0^2-\Omega^2)^2=12\beta^2\Omega^2.
\]
两个正根 $\omega_1,\omega_2$ 满足
$\omega_1\omega_2=\omega_0^2$、
$|\omega_2-\omega_1|=2\sqrt3\,\beta$。
因此
\[
\boxed{\Omega_{\rm vel,res}=\omega_0
=\sqrt{\omega_1\omega_2}},
\qquad
\boxed{\beta=\frac{|\omega_2-\omega_1|}{2\sqrt3}},
\]
\[
\boxed{\omega_d=
\sqrt{\omega_1\omega_2-
\frac{(\omega_2-\omega_1)^2}{12}}}.
\]
\excheck 需区分\emph{位移共振}与\emph{速度共振}：后者精确位于 $\omega_0$。

\oscproblem{88}{166}
某机械振动系统的位移共振曲线对应对数减量
$\lambda=1.60$。
求最大稳态位移幅度与极低驱动频率下位移幅度的比值。
\exsolution
设 $\omega_0$ 为固有频率，阻尼 $\beta$；
记 $\delta=\beta/\omega_d=\lambda/(2\pi)$。
低频幅度
$A(0)=F_0/(m\omega_0^2)$。
位移共振最大幅度
$A_{\max}=F_0/[2m\beta\sqrt{\omega_0^2-\beta^2}]$。
由 $\omega_d^2=\omega_0^2-\beta^2$，得
\[
\frac{A_{\max}}{A(0)}
=\frac{\omega_0^2}{2\beta\omega_d}
=\frac{1+\delta^2}{2\delta}.
\]
故
\[
\boxed{\frac{A_{\max}}{A(0)}
=\frac{\pi}{\lambda}+\frac{\lambda}{4\pi}
\simeq2.1}.
\]
\excheck $\lambda=1.60<2\pi$，非零位移共振峰存在。

\oscproblem{89}{166}
悬挂的弹簧振子受外力
$F_x=F_0\cos\omega t$，
稳态位移
$x=a\cos(\omega t-\varphi)$。
求外力在一个完整振动周期内做功。
\exsolution
瞬时功率 $P=F_x\dot x$，
且 $\dot x=-a\omega\sin(\omega t-\varphi)$。
令 $\theta=\omega t$，
\[
W=\int_0^{2\pi/\omega}F_x\dot x\,\mathrm dt
=-aF_0\int_0^{2\pi}\cos\theta\,
\sin(\theta-\varphi)\,\mathrm d\theta.
\]
展开正弦，周期平均
$\int_0^{2\pi}\cos^2\theta\,\mathrm d\theta=\pi$、
$\int_0^{2\pi}\sin\theta\cos\theta\,\mathrm d\theta=0$，
所以
\[
\boxed{W=\pi aF_0\sin\varphi}.
\]
\excheck 稳态周期中储能净变化为零，正功恰抵消阻尼耗能；保守力无此净输入。

\oscproblem{90}{166}
质量 $m=50\,\mathrm g$ 的球悬于弹簧
$k=20.0\,\mathrm{N/m}$，
在频率 $\omega=25.0\,\mathrm{s^{-1}}$ 的竖直周期外力作用下，
稳态位移幅度 $a=1.3\,\mathrm{cm}$，
位移相对外力滞后相位 $\varphi=3\pi/4$。
求：(a) 品质因数；(b) 外力在一个周期内做功。
\exsolution
固有频率平方
$\omega_0^2=k/m=400\,\mathrm{s^{-2}}$，
阻尼力模型为 $2m\beta\dot x$。
相位关系
$\tan\varphi=2\beta\omega/(\omega_0^2-\omega^2)$；
$\tan(3\pi/4)=-1$，
$\omega_0^2-\omega^2=-225$，得
$2\beta\omega=225$，即
$\beta=4.5\,\mathrm{s^{-1}}$。
品质因数按 $Q=\omega_0/(2\beta)$：
\[
\boxed{Q=\frac{20}{9}\simeq2.2}.
\]
由实部响应
$F_0\cos\varphi=m a(\omega_0^2-\omega^2)$，
并用第89题功公式
$W=\pi aF_0\sin\varphi$，
\[
W=\pi m a^2(\omega_0^2-\omega^2)\tan\varphi,
\quad
\boxed{W\simeq6.0\,\mathrm{mJ}}.
\]
\excheck $\varphi$ 位于第二象限，表示驱动频率高于固有频率；尽管 $\cos\varphi<0$，做功为正。

\oscproblem{91}{166}
质量为 $m$ 的阻尼弹簧振子固有角频率 $\omega_0$、
阻尼系数 $\beta$，受竖直外力
$F=F_0\cos\omega t$ 而稳态振动。
求：(a) 外力一周期平均功率；
(b) 平均功率最大的驱动频率及其值。
\exsolution
稳态位移幅度
$A=(F_0/m)/\sqrt{D}$，
$D=(\omega_0^2-\omega^2)^2+4\beta^2\omega^2$。
相位滞后满足
$\sin\varphi=2\beta\omega/\sqrt D$。
第89题每周做功为 $\pi AF_0\sin\varphi$，
除以周期 $2\pi/\omega$：
\[
\boxed{\overline P=
\frac{F_0^2\beta\omega^2/m}
{(\omega_0^2-\omega^2)^2+4\beta^2\omega^2}}.
\]
将分母除以 $\omega^2$，
可见当 $\omega=\omega_0$ 时取最小，
故
\[
\boxed{\omega_{\rm Pmax}=\omega_0,\qquad
\overline P_{\max}=\frac{F_0^2}{4m\beta}}.
\]
\excheck 平均\emph{吸收功率}共振发生在 $\omega_0$，而位移响应共振常在更低的频率。

\oscproblem{92}{166}
外力幅度保持恒定而频率可调，竖直作用于阻尼弹簧球。阻尼系数 $\beta=\omega_0/\eta$，即比固有角频率小 $\eta$ 倍。求在\emph{位移共振}频率下平均吸收功率较最大平均功率降低百分之多少。
\exsolution
位移共振频率满足
$\omega_r^2=\omega_0^2-2\beta^2$，
在 $\eta>\sqrt2$ 时存在非零峰。
由第91题
\[
\frac{\overline P(\omega_r)}{\overline P_{\max}}
=\frac{4\beta^2\omega_r^2}
{(\omega_0^2-\omega_r^2)^2+4\beta^2\omega_r^2}
=\frac{\omega_0^2-2\beta^2}{\omega_0^2-\beta^2}.
\]
所以相对差百分率
\[
\boxed{
100\frac{\overline P_{\max}-\overline P(\omega_r)}
{\overline P_{\max}}
=\frac{100}{\eta^2-1}\,\%}.
\]
\excheck 若 $\eta\le\sqrt2$ 则位移共振峰移至零频极限，须单独讨论而不直接套入此式。

\oscproblem{93}{167}
均匀水平圆盘由中心弹性竖杆悬挂，
外施谐波扭矩 $N_\tau=N_m\cos\omega t$，
扭转稳态响应
$\varphi(t)=\varphi_m\cos(\omega t-\alpha)$。
求：(a) 摩擦力在一周期内所做功；
(b) 已知圆盘对轴转动惯量为 $J$，求系统品质因数。
\exsolution
(a) 外扭矩每周做功
$W_{\rm ext}=\pi N_m\varphi_m\sin\alpha$，
稳态系统能量周期变化为零，
因此摩擦做功是其负值：
\[
\boxed{W_{\rm fric}=-\pi N_m\varphi_m\sin\alpha}.
\]
(b) 稳态旋转方程
$J\ddot\varphi+2J\beta\dot\varphi+
J\omega_0^2\varphi=N_m\cos\omega t$。
由相量实部和虚部
\[
\omega_0^2-\omega^2
=\frac{N_m\cos\alpha}{J\varphi_m},
\qquad
2\beta\omega=\frac{N_m\sin\alpha}{J\varphi_m}.
\]
依 $Q=\omega_0/(2\beta)$，
\[
\boxed{Q=\frac{\omega J\varphi_m}{N_m\sin\alpha}
\sqrt{\omega^2+\frac{N_m\cos\alpha}{J\varphi_m}}}.
\]
\excheck 只使用外力扭矩、响应幅相和惯量即可确定阻尼与固有频率；$\sin\alpha>0$ 保障耗散功为负。
% === END 23-mechanical-oscillations-082-093.tex ===
